
%\newpage
\chapter{约伯记}
\section{受苦 1-2:10}
\subsection{约伯的简介 1:1-5}
\begin{table}[H]
\hspace{-5mm}
\centering
\begin{tabular}{p{7.em}|p{13.5em}|p{28em}}    \toprule
\multicolumn{1}{c|}{品格}&\multicolumn{1}{c|}{家业}&\multicolumn{1}{c}{常為儿女自潔}\\ \midrule
烏斯地有一個人名叫約伯。那人完全、正直，敬畏神，遠離惡事。
&他生了七個兒子，三個女兒。他的家產有七千羊，三千駱駝，五百對牛，五百母驢，並有許多僕婢。這人在東方人中就為至大。
&他的兒子按著日子各在自己家裡設擺筵宴，就打發人去，請了他們的三個姐妹來，與他們一同吃喝。筵宴的日子過了，約伯打發人去叫他們自潔。他清早起來，按著他們眾人的數目獻燔祭，因為他說：「恐怕我兒子犯了罪，心中棄掉神。」約伯常常這樣行。\\    \bottomrule
\end{tabular}
\end{table}


\subsection{一次受苦 1:6-22}
\subsubsection{撒旦激動神攻擊约伯 1:6-12}
\begin{table}[H]
%\hspace{-3.8em}
\centering
\begin{tabular}{p{30.8em}|p{15.2em}}    \toprule
\multicolumn{1}{c|}{撒旦}&\multicolumn{1}{c}{耶和華}\\ \midrule
有一天，神的眾子來侍立在耶和華面前，撒旦也來在其中。 
&耶和華問撒旦說：「你從哪裡來？」\\ \midrule
撒旦回答說：「我從地上走來走去，往返而來。」&耶和華問撒旦說：「你曾用心察看我的僕人約伯沒有？地上再沒有人像他完全、正直，敬畏神，遠離惡事。」\\ \midrule
撒旦回答耶和華說：「約伯敬畏神，豈是無故呢？你豈不是四面圈上籬笆，圍護他和他的家並他一切所有的嗎？他手所做的都蒙你賜福，他的家產也在地上增多。你且伸手毀他一切所有的，他必當面棄掉你。」&耶和華對撒旦說：「凡他所有的都在你手中，只是不可伸手加害於他。」\\ \midrule
於是撒旦從耶和華面前退去。&\\    \bottomrule
\end{tabular}
\end{table}

\subsubsection{突如其来的四个灾 1:13-19}
\begin{table}[H]
\centering 
\begin{tabular}{p{13.9em}|p{8.em}|p{10.em}|p{14.4em}}    \toprule 
\multicolumn{1}{c|}{牛和驢}&\multicolumn{1}{c|}{群羊和僕人}&\multicolumn{1}{c|}{駱駝}&\multicolumn{1}{c}{兒女}\\\midrule     
有一天，約伯的兒女正在他們長兄的家裡吃飯喝酒，有報信的來見約伯說：「牛正耕地，驢在旁邊吃草，示巴人忽然闖來，把牲畜擄去，並用刀殺了僕人，唯有我一人逃脫，來報信給你。」
&他還說話的時候，又有人來說：「神從天上降下火來，將群羊和僕人都燒滅了，唯有我一人逃脫，來報信給你。」
&他還說話的時候，又有人來說：「迦勒底人分做三隊忽然闖來，把駱駝擄去，並用刀殺了僕人，唯有我一人逃脫，來報信給你。」
&他還說話的時候，又有人來說：「你的兒女正在他們長兄的家裡吃飯喝酒，不料有狂風從曠野颳來，擊打房屋的四角，房屋倒塌在少年人身上，他們就都死了，唯有我一人逃脫，來報信給你。」\\\bottomrule
\end{tabular} 
\end{table}

\subsubsection{約伯对苦难的态度 1:20-22}
約伯便起來，撕裂外袍，剃了頭，伏在地上下拜， 21 說：「我赤身出於母胎，也必赤身歸回；賞賜的是耶和華，收取的也是耶和華。耶和華的名是應當稱頌的！」 22 在這一切的事上約伯並不犯罪，也不以神為愚妄(也不妄評神)。

\subsection{二次受苦 2:1-10}
\subsubsection{撒旦二次激動神攻擊約伯 2:1-6}
\begin{table}[H]
%\hspace{-3.8em}
\centering
\begin{tabular}{p{26.8em}|p{24.em}}    \toprule
\multicolumn{1}{c|}{撒旦}&\multicolumn{1}{c}{耶和華}\\ \midrule
又有一天，神的眾子來侍立在耶和華面前，撒旦也來在其中。 
&耶和華問撒旦說：「你從哪裡來？」\\ \midrule
撒旦回答說：「我從地上走來走去，往返而來。」&耶和華問撒旦說：「你曾用心察看我的僕人約伯沒有？地上再沒有人像他完全、正直，敬畏神，遠離惡事。你雖激動我攻擊他，無故地毀滅他，他仍然持守他的純正。」\\ \midrule
撒旦回答耶和華說：「人以皮代皮，情願捨去一切所有的，保全性命。你且伸手傷他的骨頭和他的肉，他必當面棄掉你。」&耶和華對撒旦說：「他在你手中，只要存留他的性命。」\\    \bottomrule
\end{tabular}
\end{table}


\subsubsection{二次受苦 1:7-8}
於是撒旦從耶和華面前退去，擊打約伯，使他從腳掌到頭頂長毒瘡。 8 約伯就坐在爐灰中，拿瓦片刮身體。 


\subsubsection{约伯对苦难的态度 1:9-10}
他的妻子對他說：「你仍然持守你的純正嗎？你棄掉神，死了吧！」 10 約伯卻對她說：「你說話像愚頑的婦人一樣。哎！難道我們從神手裡得福，不也受禍嗎？」在這一切的事上，約伯並不以口犯罪。

\section{与三个朋友辩论 2:11-31}
\subsection{三个朋友看约伯 2:11-13}
\begin{itemize}
\itemsep-.45em  
\item 三个朋友看约伯
\begin{itemize}
\itemsep-.35em  
\item 提幔人以利法
\item 書亞人比勒達
\item 拿瑪人瑣法
\end{itemize}
\item 朋友的反应
\begin{itemize}
\itemsep-.35em  
\item 他們遠遠的舉目觀看，認不出他來，就放聲大哭
\item 各人撕裂外袍，把塵土向天揚起來，落在自己的頭上
\item 他們就同他七天七夜坐在地上，一個人也不向他說句話
\end{itemize}
\end{itemize}
約伯的三個朋友提幔人以利法、書亞人比勒達、拿瑪人瑣法，聽說有這一切的災禍臨到他身上，各人就從本處約會同來，為他悲傷，安慰他。 12 他們遠遠地舉目觀看，認不出他來，就放聲大哭。各人撕裂外袍，把塵土向天揚起來，落在自己的頭上。 13 他們就同他七天七夜坐在地上，一個人也不向他說句話，因為他極其痛苦。

\subsection{約伯恨生求死 3}
\subsubsection{恨生-11个“願”咒詛自己的生日 3:1-10}
\begin{itemize}
\itemsep-.45em  
\item 願我生的那日和說懷了男胎的那夜都滅沒
\item 願那日變為黑暗
\item 願神不從上面尋找他
\item 願亮光不照於其上
\item 願黑暗和死蔭索取那日 
\item 願密雲停在其上
\item 願日蝕恐嚇他。
\item 願那夜被幽暗奪取，不在年中的日子同樂，也不入月中的數目
\item 願那夜沒有生育，其間也沒有歡樂的聲音
\item 願那咒詛日子且能惹動鱷魚的咒詛那夜
\item 願那夜黎明的星宿變為黑暗，盼亮卻不亮，也不見早晨的光線
\end{itemize}
此後，約伯開口咒詛自己的生日， 2 說： 3 「願我生的那日和說懷了男胎的那夜都滅沒。 4 願那日變為黑暗，願神不從上面尋找它，願亮光不照於其上。 5 願黑暗和死蔭索取那日，願密雲停在其上，願日食恐嚇它。 6 願那夜被幽暗奪取，不在年中的日子同樂，也不入月中的數目。 7 願那夜沒有生育，其間也沒有歡樂的聲音。 8 願那咒詛日子且能惹動鱷魚的，咒詛那夜。 9 願那夜黎明的星宿變為黑暗，盼亮卻不亮，也不見早晨的光線 a。 10 因沒有把懷我胎的門關閉，也沒有將患難對我的眼隱藏。



\subsubsection{求死-三个“為何”  3:11-26}
\begin{itemize}
\itemsep-.45em  
\item 為何不出母胎而死？為何不出母腹絕氣？為何有膝接收我？為何有奶哺養我？ 3:11-19
\begin{itemize}
\itemsep-.45em  
\item 和地上為自己重造荒邱的君王、謀士躺臥安睡
\item 與有金子、將銀子裝滿了房屋的王子一同安息
\item 像隱而未現、不到期而落的胎，歸於無有，如同未見光的嬰孩
\item 惡人止息攪擾
\item 困乏人得享安息
\item 被囚的人同得安逸，不聽見督工的聲音
\item 大小都在那裡；奴僕脫離主人的轄制
\end{itemize}
\item 受患難的人為何有光賜給他呢？心中愁苦的人為何有生命賜給他呢？ 3:20-22
\begin{itemize}
\itemsep-.45em  
\item 他們切望死，卻不得死
\item 求死，勝於求隱藏的珍寶
\item 他們尋見墳墓就快樂，極其歡喜
\end{itemize}
\item 人的道路既然遮隱，神又把他四面圍困，為何有光賜給他呢？ 3:23-26
\begin{itemize}
\itemsep-.45em  
\item 我未曾吃飯就發出歎息；我唉哼的聲音湧出如水
\item 我所恐懼的臨到我身，我所懼怕的迎我而來
\item 我不得安逸，不得平靜，也不得安息，卻有患難來到
\end{itemize}
\end{itemize}

\begin{table}[H]
\centering 
\begin{tabular}{p{23.92em}|p{9.6em}|p{14.em}}    \toprule 
\multicolumn{1}{c|}{為何出生}&\multicolumn{1}{c|}{受患難的人為何有光}&\multicolumn{1}{c}{道路被遮隱的人為何有光}\\\midrule     
我\textcolor{red}{為何}不出母胎而死？\textcolor{red}{為何}不出母腹絕氣？\textcolor{red}{為何}有膝接收我？\textcolor{red}{為何}有奶哺養我？不然，我就早已躺臥安睡，和地上為自己重造荒丘的君王、謀士，或與有金子、將銀子裝滿了房屋的王子一同安息；或像隱而未現、不到期而落的胎，歸於無有，如同未見光的嬰孩。在那裡，惡人止息攪擾，困乏人得享安息；被囚的人同得安逸，不聽見督工的聲音。大小都在那裡，奴僕脫離主人的轄制。
&受患難的人，\textcolor{red}{為何}有光賜給他呢？心中愁苦的人，\textcolor{red}{為何}有生命賜給他呢？他們切望死卻不得死，求死勝於求隱藏的珍寶。他們尋見墳墓就快樂，極其歡喜。
&人的道路既然遮隱，神又把他四面圍困，\textcolor{red}{為何}有光賜給他呢？我未曾吃飯就發出嘆息，我唉哼的聲音湧出如水。因我所恐懼的臨到我身，我所懼怕的迎我而來。我不得安逸，不得平靜，也不得安息，卻有患難來到。\\\bottomrule
\end{tabular} 
\end{table}
 


\subsection{第一轮辯論 4-14}
\subsubsection{以利法一次回答-“神消灭耕罪孽、種毒害的人” 4-5}
\subsubsubsection{质问约伯”你过去敬畏神寻求完全，为什么祸患临到却惊恐“ 4:1-6}
提幔人以利法回答說： 2 「人若想與你說話，你就厭煩嗎？但誰能忍住不說呢？ 3 你素來教導許多的人，又堅固軟弱的手。 4 你的言語曾扶助那將要跌倒的人，你又使軟弱的膝穩固。 5 但現在禍患臨到你，你就昏迷；挨近你，你便驚惶。你的倚靠不是在你敬畏神嗎？你的盼望不是在你行事純正嗎？

\subsubsubsection{谈自己的看法和默示 4:7-21}
\begin{itemize}
\itemsep-.45em  
\item 个人经验，“無辜的人有誰滅亡”,神所行無不公義 4:7-11\\
請你追想，無辜的人有誰滅亡？正直的人在何處剪除？ 8 按我所見，耕罪孽、種毒害的人，都照樣收割。 9 神一出氣，他們就滅亡；神一發怒，他們就消沒。 10 獅子的吼叫和猛獅的聲音，盡都止息；少壯獅子的牙齒，也都敲掉。 11 老獅子因絕食而死，母獅之子也都離散。

\item 夜间的默示''根基在塵土裡、被蠹蟲所毀壞的人被毀滅，永歸無有，無人理會'' 4:12-21\\
我暗暗地得了默示，我耳朵也聽其細微的聲音。 13 在思念夜中異象之間，世人沉睡的時候， 14 恐懼、戰兢臨到我身，使我百骨打戰。 15 有靈從我面前經過，我身上的毫毛直立。 16 那靈停住，我卻不能辨其形狀，有影像在我眼前，我在靜默中聽見有聲音說： 17 『必死的人豈能比神公義嗎？人豈能比造他的主潔淨嗎？ 18 主不信靠他的臣僕，並且指他的使者為愚昧， 19 何況那住在土房，根基在塵土裡，被蠹蟲所毀壞的人呢？ 20 早晚之間，就被毀滅，永歸無有，無人理會。 21 他帳篷的繩索豈不從中抽出來呢？他死，且是無智慧而死。
\end{itemize}

\subsubsubsection{暗喻約伯患難源于愚妄,惡人終難免禍 5:1-7}
\begin{itemize}
\itemsep-.45em  
\item 忿怒害死愚妄人；嫉妒殺死癡迷人
\item 他的兒女遠離穩妥的地步，在城門口被壓，並無人搭救
\item 他的莊稼有飢餓的人吃盡了，就是在荊棘裡的也搶去了
\item 他的財寶有網羅張口吞滅了
\item 禍患原不是從土中出來；患難也不是從地裡發生
\end{itemize}
你且呼求，有誰答應你？諸聖者之中，你轉向那一位呢？ 2 憤怒害死愚妄人，嫉妒殺死痴迷人。 3 我曾見愚妄人扎下根，但我忽然咒詛他的住處。 4 他的兒女遠離穩妥的地步，在城門口被壓，並無人搭救。 5 他的莊稼有飢餓的人吃盡了，就是在荊棘裡的也搶去了，他的財寶有網羅張口吞滅了。 6 禍患原不是從土中出來，患難也不是從地裡發生。 7 人生在世必遇患難，如同火星飛騰。

\subsubsubsection{述说神的全能，以警示約伯狡詐人在神面前没有出路 5:8-16} 
至於我，我必仰望神，把我的事情託付他。 9 他行大事不可測度，行奇事不可勝數。 10 降雨在地上，賜水於田裡； 11 將卑微的安置在高處，將哀痛的舉到穩妥之地； 12 破壞狡猾人的計謀，使他們所謀的不得成就。 13 他叫有智慧的中了自己的詭計，使狡詐人的計謀速速滅亡。 14 他們白晝遇見黑暗，午間摸索如在夜間。 15 神拯救窮乏人，脫離他們口中的刀和強暴人的手。 16 這樣，貧寒的人有指望，罪孽之輩必塞口無言。

\subsubsubsection{劝約伯降伏''因犯罪''而得的管教,受神懲乃為有福 5:17-27}
神所懲治的人是有福的，所以你不可輕看全能者的管教。 18 因為他打破，又纏裹；他擊傷，用手醫治。 19 你六次遭難，他必救你，就是七次，災禍也無法害你。 20 在饑荒中，他必救你脫離死亡；在爭戰中，他必救你脫離刀劍的權力。 21 你必被隱藏，不受口舌之害；災殃臨到，你也不懼怕。 22 你遇見災害饑饉，就必喜笑；地上的野獸，你也不懼怕。 23 因為你必與田間的石頭立約，田裡的野獸也必與你和好。 24 你必知道你帳篷平安，要查看你的羊圈，一無所失。 25 也必知道你的後裔將來發達，你的子孫像地上的青草。 26 你必壽高年邁才歸墳墓，好像禾捆到時收藏。 27 這理我們已經考察，本是如此，你須要聽，要知道是與自己有益。


\subsubsection{約伯驳斥以利法 6-7（15/51节）}
\begin{itemize}
\itemsep-.45em 
\item 驳斥朋友 6:1-7:6
\begin{itemize}
\itemsep-.45em
\item 神的攻擊身心灵受到惊吓 6:1-4 \\
約伯回答說： 2 「唯願我的煩惱稱一稱，我一切的災害放在天平裡！ 3 現今都比海沙更重，所以我的言語急躁。 4 因全能者的箭射入我身，其毒，我的靈喝盡了；神的驚嚇擺陣攻擊我。
\item 朋友的安慰如可厭的食物 6:5-7\\
野驢有草豈能叫喚？牛有料豈能吼叫？ 6 物淡而無鹽豈可吃嗎？蛋青有什麼滋味呢？ 7 看為可厭的食物，我心不肯挨近。
\item 表达自己向神的两个愿望 6:8-13\\
唯願我得著所求的，願神賜我所切望的， 9 就是願神把我壓碎，伸手將我剪除！ 10 我因沒有違棄那聖者的言語，就仍以此為安慰，在不止息的痛苦中還可踴躍。 11 我有什麼氣力使我等候？我有什麼結局使我忍耐？ 12 我的氣力豈是石頭的氣力？我的肉身豈是銅的呢？ 13 在我豈不是毫無幫助嗎？智慧豈不是從我心中趕出淨盡嗎？
\item 责备朋友用诡诈的言语枉谈约伯的苦情 6:14-23\\
那將要灰心，離棄全能者，不敬畏神的人，他的朋友當以慈愛待他。 15 我的弟兄詭詐，好像溪水，又像溪水流乾的河道。這河因結冰發黑，有雪藏在其中。 17 天氣漸暖就隨時消化，日頭炎熱便從原處乾涸。 18 結伴的客旅離棄大道，順河偏行，到荒野之地死亡。 19 提瑪結伴的客旅瞻望，示巴同夥的人等候。 20 他們因失了盼望就抱愧，來到那裡便蒙羞。 21 現在你們正是這樣，看見驚嚇的事便懼怕。 22 我豈說『請你們供給我，從你們的財物中送禮物給我』？ 23 豈說『拯救我脫離敵人的手』嗎？『救贖我脫離強暴人的手』嗎？
\item 劝朋友用公正的态度辨明约伯的苦情 6:24-30\\
請你們教導我，我便不作聲，使我明白在何事上有錯。 25 正直的言語力量何其大！但你們責備是責備什麼呢？ 26 絕望人的講論既然如風，你們還想要駁正言語嗎？ 27 你們想為孤兒拈鬮，以朋友當貨物。 28 現在請你們看看我，我決不當面說謊。 29 請你們轉意，不要不公；請再轉意，我的事有理。 30 我的舌上豈有不義嗎？我的口裡豈不辨奸惡嗎？
\end{itemize}
\item 向神祷告 7:7-21 \underline{（15/51节）}
\begin{itemize}
\itemsep-.45em 
\item 述说自己的苦情 7:1-6\\
人在世上豈無爭戰嗎？他的日子不像雇工人的日子嗎？ 2 像奴僕切慕黑影，像雇工人盼望工價， 3 我也照樣經過困苦的日月，夜間的疲乏為我而定。 4 我躺臥的時候便說：『我何時起來，黑夜就過去呢？』我盡是翻來覆去，直到天亮。 5 我的肉體以蟲子和塵土為衣，我的皮膚才收了口又重新破裂。 6 我的日子比梭更快，都消耗在無指望之中。 7 求你想念我的生命不過是一口氣！我的眼睛必不再見福樂。
\item 求神想念自己生命短暂如云彩消散 7:7-11\\
求你想念我的生命不過是一口氣！我的眼睛必不再見福樂。 8 觀看我的人，他的眼必不再見我；你的眼目要看我，我卻不在了。 9 雲彩消散而過，照樣，人下陰間也不再上來。 10 他不再回自己的家，故土也不再認識他。我不禁止我口，我靈愁苦要發出言語，我心苦惱要吐露哀情。
\item 如果自己有罪，求神赦免而不是以人为箭靶子时刻监察并恐吓 7:12-21\\
我對神說：我豈是洋海，豈是大魚，你竟防守我呢？ 13 若說我的床必安慰我，我的榻必解釋我的苦情， 14 你就用夢驚駭我，用異象恐嚇我。 15 甚至我寧肯噎死，寧肯死亡，勝似留我這一身的骨頭。 16 我厭棄性命，不願永活。你任憑我吧！因我的日子都是虛空。 17 人算什麼，你竟看他為大，將他放在心上， 18 每早鑒察他，時刻試驗他？ 19 你到何時才轉眼不看我，才任憑我咽下唾沫呢？ 20 鑒察人的主啊，我若有罪，於你何妨？為何以我當你的箭靶子，使我厭棄自己的性命？ 21 為何不赦免我的過犯，除掉我的罪孽？我現今要躺臥在塵土中，你要殷勤地尋找我，我卻不在了。
\end{itemize}
\end{itemize}


\subsubsection{比勒達一次回答-或者约伯兒女得罪了神受報應 8}
\begin{itemize}
\itemsep-.45em  
\item 责备约伯言语如狂风 8:1-3\\
書亞人比勒達回答說： 2 「這些話你要說到幾時？口中的言語如狂風要到幾時呢？ 3 神豈能偏離公平？全能者豈能偏離公義？
\item 猜测约伯的祸源 8:4\\
或者你的兒女得罪了他，他使他們受報應。
\item 劝约伯殷勤寻求清洁正直 8:5-7\\
你若殷勤地尋求神，向全能者懇求， 6 你若清潔正直，他必定為你起來，使你公義的居所興旺。 7 你起初雖然微小，終久必甚發達。
\item 以古人的经验教导约伯 8:8-10\\
請你考問前代，追念他們的列祖所查究的。 9 （我們不過從昨日才有，一無所知，我們在世的日子好像影兒。） 10 他們豈不指教你，告訴你，從心裡發出言語來呢？
\item 以植物喻恶人必然被剪除 8:11-12
\begin{itemize}
\itemsep-.45em  
\item 蒲草沒有泥豈能發長
\item 蘆荻沒有水豈能生發,尚青的時候，還沒有割下，比百樣的草先枯槁
\end{itemize}
\item 忘記神,不虔敬人的结局 8:13-19\\
凡忘記神的人，景況也是這樣。不虔敬人的指望要滅沒， 14 他所仰賴的必折斷，他所倚靠的是蜘蛛網。 15 他要倚靠房屋，房屋卻站立不住；他要抓住房屋，房屋卻不能存留。 16 他在日光之下發青，蔓子爬滿了園子； 17 他的根盤繞石堆，扎入石地。 18 他若從本地被拔出，那地就不認識他，說：『我沒有見過你。』 19 看哪，這就是他道中之樂，以後必另有人從地而生。
\item 如果约伯肯作完全人的结局 8:20-22\\
神必不丟棄完全人，也不扶助邪惡人。 21 他還要以喜笑充滿你的口，以歡呼充滿你的嘴。 22 恨惡你的要披戴慚愧，惡人的帳篷必歸於無有。
\end{itemize}


\subsubsection{约伯驳斥比勒達 9-10（33/57节）}
\begin{itemize}
\itemsep-.45em
\item 约伯为自己的辩解 9:1-24
\begin{itemize}
\itemsep-.45em
\item 谁能有机会到神面前争辩自己的义 9:1-3\\
約伯回答說： 2 「我真知道是這樣，但人在神面前怎能成為義呢？ 3 若願意與他爭辯，千中之一也不能回答。
\item 述说神大能创造的作为（自然界） 9:4-10\\
他心裡有智慧，且大有能力，誰向神剛硬而得亨通呢？ 5 他發怒把山翻倒挪移，山並不知覺。 6 他使地震動離其本位，地的柱子就搖撼。 7 他吩咐日頭不出來，就不出來，又封閉眾星。 8 他獨自鋪張蒼天，步行在海浪之上。 9 他造北斗、參星、昴星，並南方的密宮。 10 他行大事不可測度，行奇事不可勝數。
\item 神对每个生命审判的主权 9:11-19\\
他從我旁邊經過，我卻不看見；他在我面前行走，我倒不知覺。 12 他奪取，誰能阻擋？誰敢問他『你做什麼』？神必不收回他的怒氣，扶助拉哈伯的屈身在他以下。 14 既是這樣，我怎敢回答他，怎敢選擇言語與他辯論呢？ 15 我雖有義，也不回答他，只要向那審判我的懇求。 16 我若呼籲，他應允我，我仍不信他真聽我的聲音。 17 他用暴風折斷我，無故地加增我的損傷。 18 我就是喘一口氣，他都不容，倒使我滿心苦惱。 19 若論力量，他真有能力；若論審判，他說：『誰能將我傳來呢？』
\item 无论善恶，神都审判 9:20-24\\
我雖有義，自己的口要定我為有罪；我雖完全，我口必顯我為彎曲。 21 我本完全，不顧自己，我厭惡我的性命。 22 善惡無分，都是一樣，所以我說，完全人和惡人，他都滅絕。 23 若忽然遭殺害之禍，他必戲笑無辜的人遇難。 24 世界交在惡人手中，蒙蔽世界審判官的臉，若不是他，是誰呢？
\end{itemize}
\item 向神祷告 9:25-10:22 \underline{（33/57节）}
\begin{itemize}
\itemsep-.45em
\item 愿神挪去使人惊恐的杖，我就不在惧怕 9:25-35\\
我的日子比跑信的更快，急速過去，不見福樂。 26 我的日子過去如快船，如急落抓食的鷹。 27 我若說我要忘記我的哀情，除去我的愁容，心中暢快， 28 我因愁苦而懼怕，知道你必不以我為無辜。 29 我必被你定為有罪，我何必徒然勞苦呢？ 30 我若用雪水洗身，用鹼潔淨我的手， 31 你還要扔我在坑裡，我的衣服都憎惡我。 32 他本不像我是人，使我可以回答他，又使我們可以同聽審判。 33 我們中間沒有聽訟的人，可以向我們兩造按手。 34 願他把杖離開我，不使驚惶威嚇我， 35 我就說話，也不懼怕他，現在我卻不是那樣。
\item 求神指示爭辯（降灾)的原因 10:1-7\\
我厭煩我的性命，必由著自己述說我的哀情，因心裡苦惱，我要說話。 2 對神說：不要定我有罪，要指示我，你為何與我爭辯。 3 你手所造的，你又欺壓，又藐視，卻光照惡人的計謀，這事你以為美嗎？ 4 你的眼豈是肉眼？你查看豈像人查看嗎？ 5 你的日子豈像人的日子？你的年歲豈像人的年歲？ 6 就追問我的罪孽，尋察我的罪過嗎？ 7 其實你知道我沒有罪惡，並沒有能救我脫離你手的。
\item 求神纪念曾经以何等的慈愛造我 10:8-12\\
你的手創造我，造就我的四肢百體，你還要毀滅我！ 9 求你記念製造我如摶泥一般，你還要使我歸於塵土嗎？ 10 你不是倒出我來好像奶，使我凝結如同奶餅嗎？ 11 你以皮和肉為衣給我穿上，用骨與筋把我全體聯絡。 12 你將生命和慈愛賜給我，你也眷顧保全我的心靈。
\item 难道神造我的目的只是为了查看我的罪孽而惩罚我吗？ 10:13-22\\
然而，你待我的這些事早已藏在你心裡，我知道你久有此意。 14 我若犯罪，你就察看我，並不赦免我的罪孽。 15 我若行惡，便有了禍；我若為義，也不敢抬頭，正是滿心羞愧，眼見我的苦情。 16 我若昂首自得，你就追捕我如獅子，又在我身上顯出奇能。 17 你重立見證攻擊我，向我加增惱怒，如軍兵更換著攻擊我。

18 「你為何使我出母胎呢？不如我當時氣絕，無人得見我。 19 這樣，就如沒有我一般，一出母胎就被送入墳墓。 20 我的日子不是甚少嗎？求你停手寬容我，叫我在往而不返之先， 21 就是往黑暗和死蔭之地以先，可以稍得暢快。 22 那地甚是幽暗，是死蔭混沌之地，那裡的光好像幽暗。
\end{itemize}
\end{itemize}

\subsubsection{瑣法一次回答 11}
\begin{itemize}
\itemsep-.45em 
\item 責約伯言語誇大 11:1-6\\
拿瑪人瑣法回答說： 2 「這許多的言語豈不該回答嗎？多嘴多舌的人豈可稱為義嗎？ 3 你誇大的話，豈能使人不作聲嗎？你戲笑的時候，豈沒有人叫你害羞嗎？ 4 你說：『我的道理純全，我在你眼前潔淨。』 5 唯願神說話，願他開口攻擊你， 6 並將智慧的奧祕指示你，他有諸般的智識。所以當知道神追討你，比你罪孽該得的還少。
\item 神智慧高深誰能測透 11:7-12\\
你考察，就能測透神嗎？你豈能盡情測透全能者嗎？ 8 他的智慧高於天，你還能做什麼？深於陰間，你還能知道什麼？ 9 其量比地長，比海寬。 10 他若經過，將人拘禁，招人受審，誰能阻擋他呢？ 11 他本知道虛妄的人，人的罪孽，他雖不留意，還是無所不見。 12 空虛的人卻毫無知識，人生在世好像野驢的駒子。
\item 勸其悔改 11:13-20\\
你若將心安正，又向主舉手—— 14 你手裡若有罪孽，就當遠遠地除掉，也不容非義住在你帳篷之中—— 15 那時你必仰起臉來，毫無斑點，你也必堅固，無所懼怕。 16 你必忘記你的苦楚，就是想起也如流過去的水一樣。 17 你在世的日子要比正午更明，雖有黑暗仍像早晨。 18 你因有指望就必穩固，也必四圍巡查，坦然安息。 19 你躺臥，無人驚嚇，且有許多人向你求恩。 20 但惡人的眼目必要失明，他們無路可逃，他們的指望就是氣絕。
%\item 约伯悔改的结局 11:15-19\\
%\item 但惡人的眼目必要失明。無路可逃；指望就是氣絕 11:20\\
\end{itemize}

\subsubsection{约伯驳斥瑣法 12-14（41/75节）}
\begin{itemize}
\itemsep-.45em 
\item 約伯自言己智不亞於友(指责朋友缪论神的作为) 12:1-13:19
\begin{itemize}
\itemsep-.45em 
\item 朋友自以为有智慧，藐視遭災的約伯 12:1-6\\
約伯回答說： 2 「你們真是子民哪！你們死亡，智慧也就滅沒了。 3 但我也有聰明，與你們一樣，並非不及你們。你們所說的，誰不知道呢？ 4 我這求告神，蒙他應允的人，竟成了朋友所譏笑的。公義完全人，竟受了人的譏笑。 5 安逸的人心裡藐視災禍，這災禍常常等待滑腳的人。 6 強盜的帳篷興旺，惹神的人穩固，神多將財物送到他們手中。
\item 走獸都能講論這一切災禍是耶和華的手 12:7-12\\
你且問走獸，走獸必指教你；又問空中的飛鳥，飛鳥必告訴你； 8 或與地說話，地必指教你；海中的魚也必向你說明。 9 看這一切，誰不知道是耶和華的手做成的呢？ 10 凡活物的生命和人類的氣息，都在他手中。 11 耳朵豈不試驗言語，正如上膛嘗食物嗎？ 12 年老的有智慧，壽高的有知識。
\item 陳述神的智慧與能力 12:13-25\\
在神有智慧和能力，他有謀略和知識。 14 他拆毀的，就不能再建造；他捆住人，便不得開釋。 15 他把水留住，水便枯乾；他再發出水來，水就翻地。 16 在他有能力和智慧，被誘惑的與誘惑人的，都是屬他。 17 他把謀士剝衣擄去，又使審判官變成愚人。 18 他放鬆君王的綁，又用帶子捆他們的腰。 19 他把祭司剝衣擄去，又使有能的人傾敗。 20 他廢去忠信人的講論，又奪去老人的聰明。 21 他使君王蒙羞被辱，放鬆有力之人的腰帶。 22 他將深奧的事從黑暗中彰顯，使死蔭顯為光明。 23 他使邦國興旺而又毀滅，他使邦國開廣而又擄去。 24 他將地上民中首領的聰明奪去，使他們在荒廢無路之地漂流。 25 他們無光，在黑暗中摸索，又使他們東倒西歪，像醉酒的人一樣。
\item 朋友編造謊言如无用的医生，闭口不言都是智慧 13:1-5\\
這一切，我眼都見過，我耳都聽過，而且明白。 2 你們所知道的，我也知道，並非不及你們。3 「我真要對全能者說話，我願與神理論。 4 你們是編造謊言的，都是無用的醫生。 5 唯願你們全然不作聲，這就算為你們的智慧！
\item 責友用炉灰的箴言妄證神為義 13:6-12\\
請你們聽我的辯論，留心聽我口中的分訴。 7 你們要為神說不義的話嗎？為他說詭詐的言語嗎？ 8 你們要為神徇情嗎？要為他爭論嗎？ 9 他查出你們來，這豈是好嗎？人欺哄人，你們也要照樣欺哄他嗎？ 10 你們若暗中徇情，他必要責備你們。 11 他的尊榮豈不叫你們懼怕嗎？他的驚嚇豈不臨到你們嗎？ 12 你們以為可記念的箴言是爐灰的箴言，你們以為可靠的堅壘是淤泥的堅壘。
\item 約伯持定自己有義願以死相争 13:13-19\\
13 「你們不要作聲，任憑我吧！讓我說話，無論如何我都承當。 14 我何必把我的肉掛在牙上，將我的命放在手中？ 15 他必殺我，我雖無指望，然而我在他面前還要辯明我所行的。 16 這要成為我的拯救，因為不虔誠的人，不得到他面前。 17 你們要細聽我的言語，使我所辯論的入你們的耳中。 18 我已陳明我的案，知道自己有義。 19 有誰與我爭論，我就情願緘默不言，氣絕而亡。
\end{itemize}
\item 向神祷告 13:20-14:22 \underline{（41/75节）}
\begin{itemize}
\itemsep-.45em 
\item 兩件不要,三个不明白 13:20-28

\begin{table}[H]
\centering 
\begin{tabular}{p{17.em}|p{32.em}}    \toprule
\multicolumn{1}{c|}{兩件不要}&\multicolumn{1}{c}{三个不明白}\\\hline
唯有兩件不要向我施行，我就不躲開你的面，就是把你的手縮回，遠離我身，又不使你的驚惶威嚇我。這樣，你呼叫，我就回答；或是讓我說話，你回答我。
&\textcolor{red}{我的罪孽和罪過有多少}呢？求你叫我知道我的過犯與罪愆。\textcolor{red}{你為何掩面}，拿我當仇敵呢？你要驚動被風吹的葉子嗎？要追趕枯乾的碎秸嗎？\textcolor{red}{你按罪狀刑罰我}，又使我擔當幼年的罪孽；也把我的腳上了木狗，並窺察我一切的道路，為我的腳掌劃定界限。我已經像滅絕的爛物，像蟲蛀的衣裳。\\ \bottomrule
\end{tabular} 
\end{table}

\item 向神的申诉人生是短暫且可悲 14:1-22\\
\begin{table}[H]
\centering 
\begin{tabular}{p{12.8em}|p{12.em}|p{12.5em}|p{10.em}}    \toprule
\multicolumn{1}{c|}{人生短少求你不看}&\multicolumn{1}{c|}{人死而无望}&\multicolumn{1}{c|}{不要纪念我的过犯}&\multicolumn{1}{c}{不要滅絕人的指望}\\\hline
人為婦人所生，日子短少，多有患難。出來如花，又被割下，飛去如影，不能存留。這樣的人，你豈睜眼看他嗎？又叫我來受審嗎？誰能使潔淨之物出於汙穢之中呢？無論誰也不能。人的日子既然限定，他的月數在你那裡，你也派定他的界限，使他不能越過，便求你轉眼不看他，使他得歇息，直等他像雇工人完畢他的日子。
&樹若被砍下，還可指望發芽，嫩枝生長不息。其根雖然衰老在地裡，幹也死在土中，及至得了水氣，還要發芽，又長枝條，像新栽的樹一樣。但人死亡而消滅，他氣絕，竟在何處呢？海中的水絕盡，江河消散乾涸。人也是如此，躺下不再起來，等到天沒有了，仍不得復醒，也不得從睡中喚醒
&唯願你把我藏在陰間，存於隱密處，等你的憤怒過去。願你為我定了日期，記念我。人若死了，豈能再活呢？我只要在我一切爭戰的日子，等我被釋放(改變)的時候來到。你呼叫，我便回答；你手所做的，你必羨慕。但如今你數點我的腳步，豈不窺察我的罪過嗎？我的過犯被你封在囊中，也縫嚴了我的罪孽。
&山崩變為無有，磐石挪開原處；水流消磨石頭，所流溢的洗去地上的塵土。你也照樣滅絕人的指望。你攻擊人常常得勝，使他去世；你改變他的容貌，叫他往而不回。他兒子得尊榮，他也不知道；降為卑，他也不覺得。但知身上疼痛，心中悲哀。\\ \bottomrule
\end{tabular} 
\end{table}


\end{itemize}
\end{itemize}

\subsection{第二轮辩论 15-21}
\subsubsection{以利法二次回答 15}
\begin{itemize}
\itemsep-.45em 
\item 再責約伯罪孽带动詭詐诈的舌頭，使自己被定罪 15:1-6\\
提幔人以利法回答說： 2 「智慧人豈可用虛空的知識回答，用東風充滿肚腹呢？ 3 他豈可用無益的話和無濟於事的言語理論呢？ 4 你是廢棄敬畏的意，在神面前阻止敬虔的心。 5 你的罪孽指教你的口，你選用詭詐人的舌頭。 6 你自己的口定你有罪，並非是我，你自己的嘴見證你的不是。
\item 质问约伯 15:7-16\\
你豈是頭一個被生的人嗎？你受造在諸山之先嗎？ 8 你曾聽見神的密旨嗎？你還將智慧獨自得盡嗎？ 9 你知道什麼是我們不知道的呢？你明白什麼是我們不明白的呢？ 10 我們這裡有白髮的和年紀老邁的，比你父親還老。 11 神用溫和的話安慰你，你以為太小嗎？ 12 你的心為何將你逼去，你的眼為何冒出火星， 13 使你的靈反對神，也任你的口發這言語？ 14 人是什麼，竟算為潔淨呢？婦人所生的是什麼，竟算為義呢？ 15 神不信靠他的眾聖者，在他眼前天也不潔淨， 16 何況那汙穢可憎、喝罪孽如水的世人呢？
\item 以列祖智慧警示约伯的灾难\underline{源于己罪} 15:17-35\\
我指示你，你要聽。我要述說所看見的， 18 就是智慧人從列祖所受、傳說而不隱瞞的。 19 （這地唯獨賜給他們，並沒有外人從他們中間經過。） 20 惡人一生之日劬勞痛苦，強暴人一生的年數也是如此。 21 驚嚇的聲音常在他耳中，在平安時，搶奪的必臨到他那裡。 22 他不信自己能從黑暗中轉回，他被刀劍等候。 23 他漂流在外求食，說：『哪裡有食物呢？』他知道黑暗的日子在他手邊預備好了。 24 急難困苦叫他害怕，而且勝了他，好像君王預備上陣一樣。 25 他伸手攻擊神，以驕傲攻擊全能者， 26 挺著頸項，用盾牌的厚凸面向全能者直闖， 27 是因他的臉蒙上脂油，腰積成肥肉。 28 他曾住在荒涼城邑，無人居住、將成亂堆的房屋。 29 他不得富足，財物不得常存，產業在地上也不加增。 30 他不得出離黑暗，火焰要將他的枝子燒乾。因神口中的氣，他要滅亡 a。 31 他不用倚靠虛假欺哄自己，因虛假必成為他的報應。 32 他的日期未到之先這事必成就，他的枝子不得青綠。 33 他必像葡萄樹的葡萄未熟而落，又像橄欖樹的花一開而謝。 34 原來不敬虔之輩必無生育，受賄賂之人的帳篷必被火燒。 35 他們所懷的是毒害，所生的是罪孽，心裡所預備的是詭詐。
\end{itemize}


\subsubsection{约伯驳斥以利法 16-17 (21/38)}
\begin{itemize}
\itemsep-.45em 
\item 責友无穷尽的虚空言語使人愁烦 16:1-5\\
約伯回答說： 2 「這樣的話我聽了許多。你們安慰人，反叫人愁煩。 3 虛空的言語有窮盡嗎？有什麼話惹動你回答呢？ 4 我也能說你們那樣的話，你們若處在我的境遇，我也會聯絡言語攻擊你們，又能向你們搖頭。 5 但我必用口堅固你們，用嘴消解你們的憂愁。
\item 歷言神使其憂苦難堪 16:6-16\\
\begin{table}[H]
\centering 
\begin{tabular}{p{5.em}|p{33.em}|p{10.em}}    \toprule
\multicolumn{1}{c|}{不能解憂}&\multicolumn{1}{c|}{神将自己作为箭靶子}&\multicolumn{1}{c}{述说自己的苦情}\\\hline
我雖說話，憂愁仍不得消解。我雖停住不說，憂愁就離開我嗎？
&但現在神使我困倦，使親友遠離我，又抓住我，做見證攻擊我。我身體的枯瘦，也當面見證我的不是。主發怒撕裂我，逼迫我，向我切齒；我的敵人怒目看我。他們向我開口，打我的臉羞辱我，聚會攻擊我。神把我交給不敬虔的人，把我扔到惡人的手中。我素來安逸，他折斷我，掐住我的頸項，把我摔碎，又立我為他的箭靶子。他的弓箭手四面圍繞我，他破裂我的肺腑，並不留情，把我的膽傾倒在地上；將我破裂又破裂，如同勇士向我直闖。&我縫麻布在我皮膚上，把我的角放在塵土中。 我的臉因哭泣發紫，在我的眼皮上有死蔭。我的手中卻無強暴，我的祈禱也是清潔。\\ \bottomrule
\end{tabular} 
\end{table}

\item 盼望与神辩论 16:17-17:9 
\begin{table}[H]
\centering 
\begin{tabular}{p{15.5em}|p{34.5em}}    \toprule
\multicolumn{1}{c|}{叹願與神辯白无路}&\multicolumn{1}{c}{申述其憂苦}\\\hline
地啊，不要遮蓋我的血！不要阻擋我的哀求！現今，在天有我的見證，在上有我的中保。我的朋友譏誚我，我卻向神眼淚汪汪。願人得與神辯白，如同人與朋友辯白一樣！因為再過幾年，我必走那往而不返之路。
&我的心靈消耗，我的日子滅盡，墳墓為我預備好了。真有戲笑我的在我這裡，我眼常見他們惹動我。願主拿憑據給我，自己為我作保；在你以外，誰肯與我擊掌呢？因你使他們心不明理，所以你必不高舉他們。控告他的朋友，以朋友為可搶奪的，連他兒女的眼睛也要失明。神使我做了民中的笑談，他們也吐唾沫在我臉上。我的眼睛因憂愁昏花，我的百體好像影兒。正直人因此必驚奇，無辜的人要興起攻擊不敬虔之輩。然而，義人要持守所行的道，手潔的人要力上加力。\\ \bottomrule
\end{tabular} 
\end{table}
\item 等候死亡來臨,看不見指望 17:10-16 \\
至於你們眾人，可以再來辯論吧！你們中間，我找不著一個智慧人。我的日子已經過了，我的謀算，我心所想望的，已經斷絕。他們以黑夜為白晝，說：『亮光近乎黑暗。』我若盼望陰間為我的房屋，若下榻在黑暗中，若對朽壞說『你是我的父』，對蟲說『你是我的母親姐妹』，這樣，我的指望在哪裡呢？我所指望的誰能看見呢？等到安息在塵土中，這指望必下到陰間的門閂那裡了。
\end{itemize}

\subsubsection{比勒達二次回答 18}
\begin{itemize}
\itemsep-.45em 
\item 斥责约伯用狂妄的言语对待朋友的劝诫 18:1-4\\
書亞人比勒達回答說： 2 「你尋索言語要到幾時呢？你可以揣摩思想，然後我們就說話。 3 我們為何算為畜生，在你眼中看做汙穢
\item 惡人的结局是滅亡 18:5-21\\
\begin{table}[H]
\centering 
\begin{tabular}{p{5.5em}|p{15.5em}|p{26.em}}    \toprule
\multicolumn{1}{c|}{亮光變黑暗}&\multicolumn{1}{c|}{腳陷入網中}&\multicolumn{1}{c}{必然滅亡}\\\hline
惡人的亮光必要熄滅，他的火焰必不照耀。他帳篷中的亮光要變為黑暗，他以上的燈也必熄滅。
&他堅強的腳步必見狹窄，自己的計謀必將他絆倒。因為他被自己的腳陷入網中，走在纏人的網羅上。圈套必抓住他的腳跟，機關必擒獲他。活扣為他藏在土內，羈絆為他藏在路上。四面的驚嚇要使他害怕，並且追趕他的腳跟。他的力量必因飢餓衰敗，禍患要在他旁邊等候。&他本身的肢體要被吞吃，死亡的長子要吞吃他的肢體。他要從所倚靠的帳篷被拔出來，帶到驚嚇的王那裡不屬他的必住在他的帳篷裡，硫磺必撒在他所住之處。下邊，他的根本要枯乾；上邊，他的枝子要剪除。他的記念在地上必然滅亡，他的名字在街上也不存留。他必從光明中被攆到黑暗裡，必被趕出世界。在本民中必無子無孫，在寄居之地也無一人存留。以後來的要驚奇他的日子，好像以前去的受了驚駭。不義之人的住處總是這樣，此乃不認識神之人的地步。\\ \bottomrule
\end{tabular} 
\end{table}
\end{itemize}


\subsubsection{約伯驳斥比勒達 19 (7/29)}
\begin{itemize}
\itemsep-.45em 
\item 責友誇己譴人 19:1-6\\
約伯回答說： 2 「你們攪擾我的心，用言語壓碎我，要到幾時呢？ 3 你們這十次羞辱我，你們苦待我也不以為恥。 4 果真我有錯，這錯乃是在我。 5 你們果然要向我誇大，以我的羞辱為證指責我， 6 就該知道是神傾覆我，用網羅圍繞我。
\item 诉说神攻击约伯并使其受苦 19:7-20
 \begin{itemize}
\itemsep-.45em 
\item 心灵经历的痛苦 19:7-12\\
我因委曲呼叫，卻不蒙應允；我呼求，卻不得公斷。 8 神用籬笆攔住我的道路，使我不得經過，又使我的路徑黑暗。 9 他剝去我的榮光，摘去我頭上的冠冕。 10 他在四圍攻擊我，我便歸於死亡，將我的指望如樹拔出來。 11 他的憤怒向我發作，以我為敵人。 12 他的軍旅一齊上來，修築戰路攻擊我，在我帳篷的四圍安營。
\item 被人离弃的痛苦 19:13-19 \\
他把我的弟兄隔在遠處，使我所認識的全然與我生疏。 14 我的親戚與我斷絕，我的密友都忘記我。 15 在我家寄居的和我的使女都以我為外人，我在他們眼中看為外邦人。 16 我呼喚僕人，雖用口求他，他還是不回答。 17 我口的氣味，我妻子厭惡；我的懇求，我同胞也憎嫌。 18 連小孩子也藐視我，我若起來，他們都嘲笑我。 19 我的密友都憎惡我，我平日所愛的人向我翻臉。
\item 身体承受的痛苦 19:20\\
我的皮肉緊貼骨頭，我只剩牙皮逃脫了。
\end{itemize}
\item 深信終得救恩 19:21-29 \\
我朋友啊，可憐我！可憐我！因為神的手攻擊我。 22 你們為什麼彷彿神逼迫我，吃我的肉還以為不足呢？唯願我的言語現在寫上，都記錄在書上， 24 用鐵筆鐫刻，用鉛灌在磐石上，直存到永遠。 25 我知道我的救贖主活著，末了必站立在地上。 26 我這皮肉滅絕之後，我必在肉體之外得見神。 27 我自己要見他，親眼要看他，並不像外人。我的心腸在我裡面消滅了！ 28 你們若說『我們逼迫他要何等地重呢？惹事的根乃在乎他』， 29 你們就當懼怕刀劍，因為憤怒惹動刀劍的刑罰，使你們知道有報應（審判）。
\end{itemize}

\subsubsection{瑣法二次回答 20}
\begin{itemize}
\itemsep-.45em 
\item 恼恨约伯的羞辱与责备 20:1-3\\
拿瑪人瑣法回答說： 2 「我心中急躁，所以我的思念叫我回答。 3 我已聽見那羞辱我責備我的話，我的悟性叫我回答。
\item 亘古以来恶人悲惨的结局 20:4-14\\
你豈不知亙古以來，自從人生在地， 5 惡人誇勝是暫時的，不敬虔人的喜樂不過轉眼之間嗎？ 6 他的尊榮雖達到天上，頭雖頂到雲中， 7 他終必滅亡，像自己的糞一樣。素來見他的人要說：『他在哪裡呢？』 8 他必飛去如夢，不再尋見；速被趕去，如夜間的異象。 9 親眼見過他的，必不再見他；他的本處，也再見不著他。 10 他的兒女要求窮人的恩，他的手要賠還不義之財。 11 他的骨頭雖然有青年之力，卻要和他一同躺臥在塵土中。

12 「他口內雖以惡為甘甜，藏在舌頭底下， 13 愛戀不捨，含在口中； 14 他的食物在肚裡卻要化為酸，在他裡面成為虺蛇的惡毒。
\item 神对待恶人的做法 20:15-29\\
他吞了財寶，還要吐出，神要從他腹中掏出來。 16 他必吸飲虺蛇的毒氣，蝮蛇的舌頭也必殺他。 17 流奶與蜜之河，他不得再見。 18 他勞碌得來的要賠還，不得享用（吞下），不能照所得的財貨歡樂。 19 他欺壓窮人，且又離棄，強取非自己所蓋的房屋（強取房屋不得再建造）。

20 「他因貪而無厭，所喜悅的連一樣也不能保守。 21 其餘的沒有一樣他不吞滅，所以他的福樂不能長久。 22 他在滿足有餘的時候，必到狹窄的地步；凡受苦楚的人，都必加手在他身上。 23 他正要充滿肚腹的時候，神必將猛烈的憤怒降在他身上；正在他吃飯的時候，要將這憤怒像雨降在他身上。 24 他要躲避鐵器，銅弓的箭要將他射透。 25 他把箭一抽，就從他身上出來，發光的箭頭從他膽中出來，有驚惶臨在他身上。 26 他的財寶歸於黑暗，人所不吹的火要把他燒滅，要把他帳篷中所剩下的燒毀。 27 天要顯明他的罪孽，地要興起攻擊他。 28 他的家產必然過去，神發怒的日子他的貨物都要消滅。 29 這是惡人從神所得的份，是神為他所定的產業。
\end{itemize}

\subsubsection{约伯驳斥瑣法 21}
\begin{itemize}
\itemsep-.45em 
\item 约伯劝朋友仔细听自己的意见 21:1-6\\
約伯回答說： 2 「你們要細聽我的言語，就算是你們安慰我。 3 請寬容我，我又要說話，說了以後，任憑你們嗤笑吧！ 4 我豈是向人訴冤，為何不焦急呢？ 5 你們要看著我而驚奇，用手摀口。 6 我每逢思想，心就驚惶，渾身戰兢。
\item 恶人的结局并非如朋友所言 21:7-18\\
惡人為何存活，享大壽數，勢力強盛呢？ 8 他們眼見兒孫，和他們一同堅立。 9 他們的家宅平安無懼，神的杖也不加在他們身上。 10 他們的公牛孳生而不斷絕，母牛下犢而不掉胎。 11 他們打發小孩子出去，多如羊群，他們的兒女踴躍跳舞。 12 他們隨著琴鼓歌唱，又因簫聲歡喜。 13 他們度日諸事亨通，轉眼下入陰間。 14 他們對神說：『離開我們吧，我們不願曉得你的道。 15 全能者是誰？我們何必侍奉他呢？求告他有什麼益處呢？』 16 看哪，他們亨通不在乎自己，惡人所謀定的離我好遠。

17 「惡人的燈何嘗熄滅？患難何嘗臨到他們呢？神何嘗發怒，向他們分散災禍呢？ 18 他們何嘗像風前的碎秸，如暴風颳去的糠秕呢？我知道你們的意思，並誣害我的計謀。 

\item 约伯对恶人的願望 21:19-27\\
你們說神為惡人的兒女積蓄罪孽，我說不如本人受報，好使他親自知道。 20 願他親眼看見自己敗亡，親自飲全能者的憤怒。 21 他的歲月既盡，他還顧他本家嗎？ 22 神既審判那在高位的，誰能將知識教訓他呢？ 23 有人至死身體強壯，盡得平靖安逸； 24 他的奶桶充滿，他的骨髓滋潤。 25 有人至死心中痛苦，終身未嘗福樂的滋味。 26 他們一樣躺臥在塵土中，都被蟲子遮蓋。
\item 约伯指责朋友设计谋陷害自己 21:28-34\\
28 你們說：『霸者的房屋在哪裡？惡人住過的帳篷在哪裡？』 29 你們豈沒有詢問過路的人嗎？不知道他們所引的證據嗎？ 30 就是惡人在禍患的日子得存留，在發怒的日子得逃脫。 31 他所行的，有誰當面給他說明？他所做的，有誰報應他呢？ 32 然而他要被抬到塋地，並有人看守墳墓。 33 他要以谷中的土塊為甘甜，在他以先去的無數，在他以後去的更多。 34 你們對答的話中既都錯謬，怎麼徒然安慰我呢？
\end{itemize}

\subsection{第三轮辩论 22-31}
\subsubsection{以利法三次回答 22}
\begin{itemize}
\itemsep-.45em 
\item 約伯因犯罪被神惩治 22:1-5\\
提幔人以利法回答說： 2 「人豈能使神有益呢？智慧人但能有益於己。 3 你為人公義，豈叫全能者喜悅呢？你行為完全，豈能使他得利呢？ 4 豈是因你敬畏他，就責備你、審判你嗎？ 5 你的罪惡豈不是大嗎？你的罪孽也沒有窮盡。
\item 歷数約伯所犯的''罪孽'' 22:6-11\\
因你無故強取弟兄的物為當頭，剝去貧寒人的衣服。 7 困乏的人，你沒有給他水喝；飢餓的人，你沒有給他食物。 8 有能力的人就得地土，尊貴的人也住在其中。 9 你打發寡婦空手回去，折斷孤兒的膀臂。 10 因此，有網羅環繞你，有恐懼忽然使你驚惶； 11 或有黑暗蒙蔽你，並有洪水淹沒你。
\item 驳斥约伯的“谬论” 22:12-20\\
神豈不是在高天嗎？你看星宿何其高呢？ 13 你說：『神知道什麼？他豈能看透幽暗施行審判呢？ 14 密雲將他遮蓋，使他不能看見；他周遊穹蒼。』 15 你要依從上古的道嗎？這道是惡人所行的。 16 他們未到死期，忽然除滅；根基毀壞，好像被江河沖去。 17 他們向神說：『離開我們吧！』又說：『全能者能把我們怎麼樣呢？』 18 哪知，神以美物充滿他們的房屋，但惡人所謀定的離我好遠。 19 義人看見他們的結局就歡喜，無辜的人嗤笑他們， 20 說：『那起來攻擊我們的果然被剪除，其餘的都被火燒滅。』
\item 勸其認識神而得平安 22:21-30\\
你要認識神，就得平安，福氣也必臨到你。 22 你當領受他口中的教訓，將他的言語存在心裡。 23 你若歸向全能者，從你帳篷中遠除不義，就必得建立。 24 要將你的珍寶丟在塵土裡，將俄斐的黃金丟在溪河石頭之間， 25 全能者就必為你的珍寶，做你的寶銀。 26 你就要以全能者為喜樂，向神仰起臉來。 27 你要禱告他，他就聽你，你也要還你的願。 28 你定意要做何事，必然給你成就，亮光也必照耀你的路。 29 人使你降卑，你仍可說必得高升；謙卑的人，神必然拯救。 30 人非無辜，神且要搭救他，他因你手中清潔，必蒙拯救。
\end{itemize}

\subsubsection{约伯驳斥以利法 23-24}
\begin{itemize}
\itemsep-.45em 
\item 自诉求告神却无门,神使得自己惊慌丧胆 23
\begin{itemize}
\itemsep-.45em 
\item 约伯相信神必理会自己的倾诉 23:1-7\\
約伯回答說： 2 「如今我的哀告還算為悖逆，我的責罰比我的唉哼還重。 3 唯願我能知道在哪裡可以尋見神，能到他的臺前， 4 我就在他面前將我的案件陳明，滿口辯白。 5 我必知道他回答我的言語，明白他向我所說的話。 6 他豈用大能與我爭辯嗎？必不這樣，他必理會我。 7 在他那裡，正直人可以與他辯論，這樣，我必永遠脫離那審判我的。
\item 自诉欲求告神却无门 23:8-9\\
只是我往前行，他不在那裡；往後退，也不能見他。 9 他在左邊行事，我卻不能看見；在右邊隱藏，我也不能見他。
\item 叙述自己对神的信靠 23:10-12
然而他知道我所行的路，他試煉我之後，我必如精金。我腳追隨他的步履，我謹守他的道，並不偏離。 12 他嘴唇的命令，我未曾背棄；我看重他口中的言語，過於我需用的飲食。 
\item 神定意降灾，使得自己惊慌丧胆 23:13-17\\
13 只是他心志已定，誰能使他轉意呢？他心裡所願的，就行出來。 14 他向我所定的，就必做成。這類的事，他還有許多。 15 所以我在他面前驚惶，我思念這事，便懼怕他。 16 神使我喪膽，全能者使我驚惶。 17 我的恐懼不是因為黑暗，也不是因為幽暗蒙蔽我的臉。
\end{itemize}
\item 祈求神刑法恶人以证明自己对神的认识是正确的 24
\begin{itemize}
\itemsep-.45em 
\item 恶人的作为 24:1-17\\
全能者既定期罰惡，為何不使認識他的人看見那日子呢？ 2 有人挪移地界，搶奪群畜而牧養。 3 他們拉去孤兒的驢，強取寡婦的牛為當頭。 4 他們使窮人離開正道，世上的貧民盡都隱藏。 5 這些貧窮人，如同野驢出到曠野，殷勤尋找食物。他們靠著野地給兒女糊口， 6 收割別人田間的禾稼，摘取惡人餘剩的葡萄。 7 終夜赤身無衣，天氣寒冷毫無遮蓋， 8 在山上被大雨淋濕，因沒有避身之處就挨近磐石。 9 又有人從母懷中搶奪孤兒，強取窮人的衣服為當頭， 10 使人赤身無衣到處流行，且因飢餓扛抬禾捆， 11 在那些人的圍牆內造油，榨酒自己還口渴。 12 在多民的城內有人唉哼，受傷的人哀號，神卻不理會那惡人的愚妄。又有人背棄光明，不認識光明的道，不住在光明的路上。 14 殺人的黎明起來，殺害困苦窮乏人，夜間又做盜賊。 15 姦夫等候黃昏，說：『必無眼能見我』，就把臉蒙蔽。 16 盜賊黑夜挖窟窿，白日躲藏，並不認識光明。 17 他們看早晨如幽暗，因為他們曉得幽暗的驚駭。
\item 惡人被高舉不過片時 24:18-25\\
這些惡人猶如浮萍快快飄去，他們所得的份在世上被咒詛，他們不得再走葡萄園的路。 19 乾旱炎熱消沒雪水，陰間也如此消沒犯罪之輩。 20 懷他的母(胎)要忘記他，蟲子要吃他覺得甘甜；他不再被人記念，不義的人必如樹折斷。 21 他惡待(他吞滅)不懷孕不生養的婦人，不善待寡婦。 22 然而神用能力保全有勢力的人，那性命難保的人仍然興起。 23 神使他們安穩，他們就有所倚靠，神的眼目也看顧他們的道路。 24 他們被高舉，不過片時就沒有了；他們降為卑，被除滅，與眾人一樣，又如穀穗被割。 25 若不是這樣，誰能證實我是說謊的，將我的言語駁為虛空呢？
\end{itemize}
\end{itemize}


\subsubsection{比勒達三次回答 25}
\begin{table}[H]
\centering 
\begin{tabular}{p{20.em}|p{29.em}}    \toprule 
\multicolumn{1}{c|}{神的威严}& \multicolumn{1}{c}{人的污秽}\\\cline{1-2}
書亞人比勒達回答說： 2 「神有治理之權，有威嚴可畏，他在高處施行和平。 3 他的諸軍，豈能數算？他的光亮一發，誰不蒙照呢
&這樣，在神面前人怎能稱義？婦人所生的怎能潔淨？ 5 在神眼前，月亮也無光亮，星宿也不清潔。 6 何況如蟲的人，如蛆的世人呢？」\\ \bottomrule
\end{tabular} 
\end{table}


\subsubsection{约伯驳斥比勒達 26-31}
\begin{itemize}
\itemsep-.45em 
\item 讽刺朋友无知 26
\begin{table}[H]
\centering 
\begin{tabular}{p{15.em}|p{35.em}}    \toprule 
\multicolumn{1}{c|}{訕笑比勒達的愚昧1-4}& \multicolumn{1}{c}{述说神的能力和作为5-14}\\\cline{1-2}
約伯回答說： 2 「無能的人蒙你何等的幫助！膀臂無力的人蒙你何等的拯救！ 3 無智慧的人蒙你何等的指教！你向他多顯大知識。 4 你向誰發出言語來？誰的靈從你而出？
&在大水和水族以下的陰魂戰兢。 6 在神面前陰間顯露，滅亡也不得遮掩。 7 神將北極鋪在空中，將大地懸在虛空。 8 將水包在密雲中，雲卻不破裂。 9 遮蔽他的寶座，將雲鋪在其上。 10 在水面的周圍劃出界限，直到光明黑暗的交界。 11 天的柱子因他的斥責，震動驚奇。 12 他以能力攪動(平靜)大海，他藉知識打傷拉哈伯； 13 藉他的靈使天有裝飾，他的手刺殺快蛇。 14 看哪，這不過是神工作的些微，我們所聽於他的是何等細微的聲音！他大能的雷聲誰能明透呢？\\ \bottomrule
\end{tabular} 
\end{table}
\item 指教朋友什么才是神的作为 27
\begin{table}[H]
\centering 
\begin{tabular}{p{12.em}|p{10.5em}|p{25.6em}}    \toprule 
\multicolumn{1}{c|}{重申神使自己受苦1-6}& \multicolumn{1}{c|}{向神的愿望7-10}&\multicolumn{1}{c}{指教朋友神的作為11-23}\\\cline{1-3}
約伯接著說：「神奪去我的理，全能者使我心中愁苦。我指著永生的神起誓（我的生命尚在我裡面，神所賜呼吸之氣仍在我的鼻孔內）：我的嘴決不說非義之言，我的舌也不說詭詐之語。我斷不以你們為是，我至死必不以自己為不正。我持定我的義，必不放鬆；在世的日子，我心必不責備我。
&願我的仇敵如惡人一樣，願那起來攻擊我的如不義之人一般。願我的仇敵如惡人一樣，願那起來攻擊我的如不義之人一般。不敬虔的人雖然得利，神奪取其命的時候，還有什麼指望呢？患難臨到他，神豈能聽他的呼求？他豈以全能者為樂，隨時求告神呢？
&神的作為，我要指教你們；全能者所行的，我也不隱瞞。你們自己也都見過，為何全然變為虛妄呢？神為惡人所定的份，強暴人從全能者所得的報(產業)乃是這樣：倘或他的兒女增多，還是被刀所殺；他的子孫必不得飽食。他所遺留的人必死而埋葬，他的寡婦也不哀哭。他雖積蓄銀子如塵沙，預備衣服如泥土，他只管預備，義人卻要穿上，他的銀子無辜的人要分取。 他建造房屋如蟲做窩，又如守望者所搭的棚。他雖富足躺臥，卻不得收殮，轉眼之間就不在了。驚恐如波濤將他追上，暴風在夜間將他颳去。東風把他飄去，又颳他離開本處。神要向他射箭，並不留情，他恨不得逃脫神的手。人要向他拍掌，並要發叱聲，使他離開本處。\\ \bottomrule
\end{tabular} 
\end{table}

\item 真正的智慧从神而出 28

\begin{table}[H]
\centering 
\begin{tabular}{p{19.em}|p{15.8em}|p{14.6em}}    \toprule 
\multicolumn{1}{c|}{人能挖宝却不尋智慧1-11}& \multicolumn{1}{c|}{不在今生精金不足與之兌換12-19}&\multicolumn{1}{c}{敬畏主遠離惡是聪明20-28}\\\cline{1-3}
銀子有礦，煉金有方。 2 鐵從地裡挖出，銅從石中熔化。 3 人為黑暗定界限，查究幽暗陰翳的石頭，直到極處。 4 在無人居住之處刨開礦穴，過路的人也想不到他們，又與人遠離，懸在空中搖來搖去。 5 至於地，能出糧食，地內好像被火翻起來。 6 地中的石頭有藍寶石，並有金沙。 7 礦中的路，鷙鳥不得知道，鷹眼也未見過； 8 狂傲的野獸未曾行過，猛烈的獅子也未曾經過。 9 人伸手鑿開堅石，傾倒山根。 10 在磐石中鑿出水道，親眼看見各樣寶物。 11 他封閉水，不得滴流，使隱藏的物顯露出來。
&然而，智慧有何處可尋？聰明之處在哪裡呢？ 13 智慧的價值無人能知，在活人之地也無處可尋。 14 深淵說：『不在我內』，滄海說：『不在我中』。智慧非用黃金可得，也不能平白銀為它的價值。 16 俄斐金和貴重的紅瑪瑙，並藍寶石，不足與較量。 17 黃金和玻璃，不足與比較；精金的器皿，不足與兌換。 18 珊瑚、水晶都不足論，智慧的價值勝過珍珠（紅寶石）。 19 古實的紅璧璽，不足與比較；精金，也不足與較量。
&智慧從何處來呢？聰明之處在哪裡呢？ 21 是向一切有生命的眼目隱藏，向空中的飛鳥掩蔽。 22 滅沒和死亡說：『我們風聞其名。』神明白智慧的道路，曉得智慧的所在。 24 因他鑒察直到地極，遍觀普天之下。 25 要為風定輕重，又度量諸水。 26 他為雨露定命令，為雷電定道路。 27 那時他看見智慧，而且述說；他堅定，並且查究。 28 他對人說：『敬畏主就是智慧，遠離惡便是聰明。』\\ \bottomrule
\end{tabular} 
\end{table}
\item 回顾自己的过去 29

\begin{table}[H]
\centering 
\begin{tabular}{p{10.em}|p{7.8em}|p{11.8em}|p{16.6em}}    \toprule 
\multicolumn{1}{c|}{有神的同在之福1-6}& \multicolumn{1}{c|}{备受尊敬7-11}&\multicolumn{1}{c|}{弱者依靠行公义12-17}&\multicolumn{1}{c}{有荣耀尊荣18-25}\\\cline{1-3}
約伯又接著說： 2 「唯願我的景況如從前的月份，如神保守我的日子。 3 那時他的燈照在我頭上，我藉他的光行過黑暗。 4 我願如壯年的時候，那時我在帳篷中，神待我有密友之情。 5 全能者仍與我同在，我的兒女都環繞我。 6 奶多可洗我的腳，磐石為我出油成河。
&我出到城門，在街上設立座位， 8 少年人見我而迴避，老年人也起身站立； 9 王子都停止說話，用手摀口； 10 首領靜默無聲，舌頭貼住上膛。 11 耳朵聽我的就稱我有福，眼睛看我的便稱讚我，
&因我拯救哀求的困苦人和無人幫助的孤兒。 13 將要滅亡的為我祝福，我也使寡婦心中歡樂。我以公義為衣服，以公平為外袍和冠冕。 15 我為瞎子的眼，瘸子的腳。 16 我為窮乏人的父，素不認識的人，我查明他的案件。 17 我打破不義之人的牙床，從他牙齒中奪了所搶的。
&我便說：『我必死在家中（窩中），必增添我的日子多如塵沙。 19 我的根長到水邊，露水終夜沾在我的枝上。 20 我的榮耀在身上增新，我的弓在手中日強。』 21 人聽見我而仰望，靜默等候我的指教。 22 我說話之後，他們就不再說，我的言語像雨露滴在他們身上。 23 他們仰望我如仰望雨，又張開口如切慕春雨。 24 他們不敢自信，我就向他們含笑；他們不使我臉上的光改變。 25 我為他們選擇道路，又坐首位；我如君王在軍隊中居住，又如弔喪的安慰傷心的人。\\ \bottomrule
\end{tabular} 
\end{table}
\item 再说现在的苦情 30:1-19
\begin{table}[H]
\centering 
\begin{tabular}{p{36.em}|p{13.6em}}    \toprule 
\multicolumn{1}{c|}{被年少卑贱人戲笑1-14}& \multicolumn{1}{c}{心極其悲傷15-19}\\\cline{1-2}
但如今，比我年少的人戲笑我，其人之父我曾藐視，不肯安在看守我羊群的狗中。 2 他們壯年的氣力既已衰敗，其手之力與我何益呢？ 3 他們因窮乏飢餓，身體枯瘦，在荒廢淒涼的幽暗中啃乾燥之地， 4 在草叢之中採鹹草，羅騰（鬆類）的根為他們的食物。 5 他們從人中被趕出，人追喊他們如賊一般， 6 以致他們住在荒谷之間，在地洞和巖穴中， 7 在草叢中叫喚，在荊棘下聚集。 8 這都是愚頑下賤人的兒女，他們被鞭打，趕出境外。現在這些人以我為歌曲，以我為笑談。 10 他們厭惡我，躲在旁邊站著，不住地吐唾沫在我臉上， 11 鬆開他們的繩索苦待我，在我面前脫去轡頭。 12 這等下流人在我右邊起來，推開我的腳，築成戰路來攻擊我。 13 這些無人幫助的，毀壞我的道，加增我的災。 14 他們來，如同闖進大破口，在毀壞之間滾在我身上。
&15 驚恐臨到我，驅逐我的尊榮如風，我的福祿如雲過去。現在我心極其悲傷，困苦的日子將我抓住。 17 夜間我裡面的骨頭刺我，疼痛不止，好像啃我。 18 因神的大力，我的外衣汙穢不堪，又如裡衣的領子將我纏住。 19 神把我扔在淤泥中，我就像塵土和爐灰一般。 
\\ \bottomrule
\end{tabular} 
\end{table}
\item 重新转向神哀求 30:20-31:40

\begin{table}[H]
\centering 
\begin{tabular}{p{11.em}|p{38.6em}}    \toprule 
\multicolumn{1}{c|}{言神不聽呼求30:20-31}& \multicolumn{1}{c}{自言雖遭困苦守道不偏31:1-40}\\\cline{1-2}
主啊，我呼求你，你不應允我；我站起來，你就定睛看我。你向我變心，待我殘忍，又用大能追逼我，把我提在風中，使我駕風而行，又使我消滅在烈風中。我知道要使我臨到死地，到那為眾生所定的陰宅。
然而，人仆倒豈不伸手，遇災難豈不求救呢？人遭難，我豈不為他哭泣呢？人窮乏，我豈不為他憂愁呢？我仰望得好處，災禍就到了；我等待光明，黑暗便來了。我心裡煩擾不安，困苦的日子臨到我身。我沒有日光就哀哭行去（我面發黑並非因日曬），我在會中站著求救。我與野狗為弟兄，與鴕鳥為同伴。我的皮膚黑而脫落，我的骨頭因熱燒焦。所以我的琴音變為悲音，我的簫聲變為哭聲。
&我與眼睛立約，怎能戀戀瞻望處女呢？ 2 從至上的神所得之份，從至高全能者所得之業，是什麼呢？ 3 豈不是禍患臨到不義的，災害臨到作孽的呢？ 4 神豈不是察看我的道路，數點我的腳步呢？
5 「我若與虛謊同行，腳若追隨詭詐， 6 我若被公道的天平稱度，使神可以知道我的純正， 7 我的腳步若偏離正路，我的心若隨著我的眼目，若有玷汙沾在我手上， 8 就願我所種的有別人吃，我田所產的被拔出來。
9 「我若受迷惑，向婦人起淫念，在鄰舍的門外蹲伏， 10 就願我的妻子給別人推磨，別人也與她同室。 11 因為這是大罪，是審判官當罰的罪孽。 12 這本是火焚燒，直到毀滅，必拔除我所有的家產。 13 我的僕婢與我爭辯的時候，我若藐視不聽他們的情節， 14 神興起，我怎樣行呢？他察問，我怎樣回答呢？ 15 造我在腹中的，不也是造他嗎？將他與我摶在腹中的，豈不是一位嗎？
16 「我若不容貧寒人得其所願，或叫寡婦眼中失望， 17 或獨自吃我一點食物，孤兒沒有與我同吃 18 （從幼年時孤兒與我同長，好像父子一樣，我從出母腹就扶助 a寡婦）， 19 我若見人因無衣死亡，或見窮乏人身無遮蓋， 20 我若不使他因我羊的毛得暖，為我祝福， 21 我若在城門口見有幫助我的，舉手攻擊孤兒， 22 情願我的肩頭從缺盆骨脫落，我的膀臂從羊矢骨折斷。 23 因神降的災禍使我恐懼，因他的威嚴我不能妄為。
24 「我若以黃金為指望，對精金說『你是我的倚靠』， 25 我若因財物豐裕，因我手多得資財而歡喜， 26 我若見太陽發光，明月行在空中， 27 心就暗暗被引誘，口便親手 28 （這也是審判官當罰的罪孽，又是我背棄在上的神）， 29 我若見恨我的遭報就歡喜，見他遭災便高興 30 （我沒有容口犯罪，咒詛他的生命）， 31 若我帳篷的人未嘗說『誰不以主人的食物吃飽呢？』 32 （從來我沒有容客旅在街上住宿，卻開門迎接行路的人）， 33 我若像亞當 b遮掩我的過犯，將罪孽藏在懷中， 34 因懼怕大眾，又因宗族藐視我使我驚恐，以致閉口無言，杜門不出—— 35 唯願有一位肯聽我！看哪，在這裡有我所畫的押，願全能者回答我！ 36 願那敵我者所寫的狀詞在我這裡，我必帶在肩上，又綁在頭上為冠冕。 37 我必向他述說我腳步的數目，必如君王進到他面前。 38 我若奪取田地，這地向我喊冤，犁溝一同哭泣， 39 我若吃地的出產不給價值，或叫原主喪命， 40 願這地長蒺藜代替麥子，長惡草代替大麥。」約伯的話說完了。 
\\ \bottomrule
\end{tabular} 
\end{table}
\end{itemize}

\section{以利戶的劝勉 32-37}
\subsection{以利戶出场 32:1-5}
於是這三個人，因約伯自以為義就不再回答他。 2 那時有布西人蘭族巴拉迦的兒子以利戶向約伯\textcolor{red}{發怒}，因約伯自以為義，不以神為義。 3 他又向約伯的三個朋友\textcolor{red}{發怒}，因為他們想不出回答的話來，仍以約伯為有罪。 4 以利戶要與約伯說話，就等候他們，因為他們比自己年老。 5 以利戶見這三個人口中無話回答，就\textcolor{red}{怒氣發作}。


\subsection{说给三个朋友的话 32:6-22}
%\subsubsection{说给三个朋友的话 32:7-33:7}
\begin{itemize}
\itemsep-.45em 
\item 我年輕，你們老邁；因此我退讓，不敢向你們陳說我的意見
\item 在人裡面有靈；全能者的氣使人有聰明
\item 尊貴的不都有智慧；壽高的不都能明白公平
\item 你們中間無一人折服約伯，駁倒他的話
\item 約伯沒有向我爭辯；我也不用你們的話回答他
\item 我不再等候要陳說我的意見。使我舒暢
\item 我必不看人的情面，也不奉承人
\end{itemize}
布西人巴拉迦的兒子以利戶回答說：「我年輕，你們老邁，因此我退讓，不敢向你們陳說我的意見。 7 我說，年老的當先說話，壽高的當以智慧教訓人。 8 但在人裡面有靈，全能者的氣使人有聰明。 9 尊貴的不都有智慧，壽高的不都能明白公平。 10 因此我說，你們要聽我言，我也要陳說我的意見。

11 「你們查究所要說的話，那時我等候你們的話，側耳聽你們的辯論， 12 留心聽你們。誰知你們中間無一人折服約伯，駁倒他的話。 13 你們切不可說：『我們尋得智慧，神能勝他，人卻不能。』 14 約伯沒有向我爭辯，我也不用你們的話回答他。

15 「他們驚奇，不再回答，一言不發。 16 我豈因他們不說話，站住不再回答，仍舊等候呢？ 17 我也要回答我的一份話，陳說我的意見。 18 因為我的言語滿懷，我裡面的靈激動我。 19 我的胸懷如盛酒之囊，沒有出氣之縫，又如新皮袋快要破裂。 20 我要說話，使我舒暢，我要開口回答。 21 我必不看人的情面，也不奉承人。 22 我不曉得奉承，若奉承，造我的主必快快除滅我。


\subsection{说给约伯的话 33-37}
\subsubsection{开场白 33:1-7}
約伯啊，請聽我的話，留心聽我一切的言語。 2 我現在開口，用舌發言。 3 我的言語要發明心中所存的正直，我所知道的，我嘴唇要誠實地說出。 4 神的靈造我，全能者的氣使我得生。 5 你若回答我，就站起來，在我面前陳明。 6 我在神面前與你一樣，也是用土造成。 7 我不用威嚴驚嚇你，也不用勢力重壓你。

\subsubsection{驳“我是清潔無過的神找機會攻擊我” 33:8-33}
\subsubsubsection{約伯的言語 33:8-12}
你所說的我聽見了，也聽見你的言語說： 9 『我是清潔無過的，我是無辜的，在我裡面也沒有罪孽。 10 神找機會攻擊我，以我為仇敵， 11 把我的腳上了木狗，窺察我一切的道路。』 12 我要回答你說，你這話無理，因神比世人更大。
\subsubsubsection{神說一次、兩次，世人卻不理會 33:13-22}
你為何與他爭論呢？因他的事都不對人解說。 14 神說一次兩次，世人卻不理會。 15 人躺在床上沉睡的時候，神就用夢和夜間的異象， 16 開通他們的耳朵，將當受的教訓印在他們心上， 17 好叫人不從自己的謀算，不行驕傲的事(將驕傲向人隱藏)， 18 攔阻人不陷於坑裡，不死在刀下。
人在床上被懲治，骨頭中不住地疼痛， 20 以致他的口厭棄食物，心厭惡美味。 21 他的肉消瘦，不得再見，先前不見的骨頭都凸出來。 22 他的靈魂臨近深坑，他的生命近於滅命的。
\begin{itemize}
\itemsep-.45em 
\item 神就用夢和夜間的異象開通他們的耳,攔阻人不陷於坑裡，不死在刀下
\item 神以肉身的疾病痛苦挽救人
\begin{itemize}
\itemsep-.45em 
\item 人在床上被懲治，骨頭中不住地疼痛
\item 口厭棄食物，心厭惡美味
\item 肉消瘦，不得再見；先前不見的骨頭都凸出來
\item 靈魂臨近深坑；他的生命近於滅命的
\end{itemize}
\end{itemize}

\subsubsubsection{天使中若有作傳話的，指示人所當行的事,神就給他開恩 23-33}
一千天使中，若有一個做傳話的與神同在，指示人所當行的事， 24 神就給他開恩，說：『救贖他免得下坑，我已經得了贖價。』 25 他的肉要比孩童的肉更嫩，他就返老還童。 26 他禱告神，神就喜悅他，使他歡呼朝見神的面，神又看他為義。 27 他在人前歌唱說：『我犯了罪，顛倒是非，這竟與我無益。 28 神救贖我的靈魂免入深坑，我的生命也必見光。』

神兩次三次向人行這一切的事， 30 為要從深坑救回人的靈魂，使他被光照耀與活人一樣。 31 約伯啊，你當側耳聽我的話，不要作聲，等我講說。 32 你若有話說，就可以回答我；你只管說，因我願以你為是。 33 若不然，你就聽我說，你不要作聲，我便將智慧教訓你。

\subsubsection{驳“我是公義，神奪去我的理” 34:1-37}

\subsubsubsection{约伯的错误言论和行为 34:1-9}
以利戶又說： 2 「你們智慧人，要聽我的話；有知識的人，要留心聽我說。 3 因為耳朵試驗話語，好像上膛嘗食物。 4 我們當選擇何為是，彼此知道何為善。 5 約伯曾說：『我是公義，神奪去我的理。 6 我雖有理，還算為說謊言的；我雖無過，受的傷還不能醫治。』 7 誰像約伯，喝譏誚如同喝水呢？ 8 他與作孽的結伴，和惡人同行。 9 他說，人以神為樂總是無益。
\subsubsubsection{神是公义的，必不作恶 34:10-12}
所以你們明理的人，要聽我的話：神斷不至行惡，全能者斷不至作孽。 11 他必按人所做的報應人，使各人照所行的得報。 12 神必不作惡，全能者也不偏離公平。
\subsubsubsection{神的主权不容挑战 34:13-20}
\begin{table}[H]
\centering
\begin{tabular}{p{1.5cm}|p{2.75cm}|p{4.00cm}|p{8.00cm}}\hline\hline
\multicolumn{1}{c|}{掌管全地}&\multicolumn{1}{c|}{掌管人生命}&\multicolumn{1}{c|}{掌管公义}&\multicolumn{1}{c}{掌管高低贵贱}\\\hline
誰派他治理地，安定全世界呢？ 
&他若專心為己，將靈和氣收歸自己， 15 凡有血氣的就必一同死亡，世人必仍歸塵土。
&你若明理，就當聽我的話，留心聽我言語的聲音。 17 難道恨惡公平的可以掌權嗎？那有公義的，有大能的，豈可定他有罪嗎？
&他對君王說：『你是鄙陋的！』對貴臣說：『你是邪惡的！』 19 他待王子不徇情面，也不看重富足的過於貧窮的，因為都是他手所造。 20 在轉眼之間，半夜之中，他們就死亡，百姓被震動而去世，有權力的被奪去非藉人手。\\\hline
\end{tabular}
\end{table} 

\subsubsubsection{神审判人的方式 34:21-30}
神注目觀看人的道路，看明人的腳步。 22 沒有黑暗、陰翳能給作孽的藏身。 23 神審判人，不必使人到他面前再三鑒察。 24 他用難測之法打破有能力的人，設立別人代替他們。 25 他原知道他們的行為，使他們在夜間傾倒滅亡。 26 他在眾人眼前擊打他們，如同擊打惡人一樣。 27 因為他們偏行不跟從他，也不留心他的道， 28 甚至使貧窮人的哀聲達到他那裡，他也聽了困苦人的哀聲。 29 他使人安靜，誰能擾亂 a呢？他掩面，誰能見他呢？無論待一國或一人，都是如此， 30 使不虔敬的人不得做王，免得有人牢籠百姓。
\subsubsubsection{规劝约伯降伏于神的主权之下 34:31-37}
有誰對神說『我受了責罰，不再犯罪。 32 我所看不明的，求你指教我。我若作了孽，必不再做』？ 33 他施行報應，豈要隨你的心願，叫你推辭不受嗎？選定的是你，不是我，你所知道的只管說吧！ 34 明理的人和聽我話的智慧人，必對我說： 35 『約伯說話沒有知識，言語中毫無智慧。』 36 願約伯被試驗到底，因他回答像惡人一樣。 37 他在罪上又加悖逆，在我們中間拍手，用許多言語輕慢神。



\subsubsection{驳“我不犯罪比犯罪有什麼好處呢” 35:1-16}
\subsubsubsection{约伯的错误 35:1-3}
以利戶又說： 2 「你以為有理，或以為你的公義勝於神的公義， 3 才說：『這與我有什麼益處？我不犯罪比犯罪有什麼好處呢？』
\subsubsubsection{人的善恶不能使神受益受损 35:4-8}
我要回答你和在你這裡的朋友。 5 你要向天觀看，瞻望那高於你的穹蒼。 6 你若犯罪，能使神受何害呢？你的過犯加增，能使神受何損呢？ 7 你若是公義，還能加增他什麼呢？他從你手裡還接受什麼呢？ 8 你的過惡，或能害你這類的人；你的公義，或能叫世人得益處。
\subsubsubsection{人被神管教时的态度 35:9-16}
\begin{table}[H]
\centering
\begin{tabular}{p{2.4cm}|p{4.500cm}|p{10.00cm}}\hline\hline
\multicolumn{1}{c|}{通常的态度求救}&\multicolumn{1}{c|}{好的态度}&\multicolumn{1}{c}{虛妄的呼求神不垂聽}\\\hline
人因多受欺壓就哀求，因受能者的轄制(膀臂)便求救，
&卻無人說：『造我的神在哪裡？他使人夜間歌唱， 11 教訓我們勝於地上的走獸，使我們有聰明勝於空中的飛鳥。』
&他們在那裡因惡人的驕傲呼求，卻無人答應。 13 虛妄的呼求，神必不垂聽，全能者也必不眷顧。 14 何況你說，你不得見他，你的案件在他面前，你等候他吧！ 15 但如今因他未曾發怒降罰，也不甚理會狂傲， 16 所以約伯開口說虛妄的話，多發無知識的言語。\\\hline
\end{tabular}
\end{table} 


\subsubsection{劝约伯“不可因贖價大就偏行不服責罰” 36:1-24}
\subsubsubsection{劝约伯再安静，因自己還有話為神說 36:1-4}
以利戶又接著說： 2 「你再容我片時，我就指示你，因我還有話為神說。 3 我要將所知道的從遠處引來，將公義歸給造我的主。 4 我的言語真不虛謊，有知識全備的與你同在。
\subsubsubsection{神為困苦人伸冤,他時常看顧義人 36:5-11}
神有大能，並不藐視人，他的智慧甚廣。 6 他不保護惡人的性命，卻為困苦人申冤。 7 他時常看顧義人，使他們和君王同坐寶座，永遠要被高舉。 8 他們若被鎖鏈捆住，被苦難的繩索纏住， 9 他就把他們的作為和過犯指示他們，叫他們知道有驕傲的行動。 10 他也開通他們的耳朵，得受教訓，吩咐他們離開罪孽轉回。 11 他們若聽從侍奉他，就必度日亨通，歷年福樂；
\subsubsubsection{人若不聽從神的管教，就要被刀殺滅，無知無識而死 36:12-14}
若不聽從，就要被刀殺滅，無知無識而死。 13 那心中不敬虔的人積蓄怒氣，神捆綁他們，他們竟不求救， 14 必在青年時死亡，與汙穢人一樣喪命。
\subsubsubsection{神对约伯的心意 36:15-24}
\begin{table}[H]
\centering
\begin{tabular}{p{4.15cm}|p{8.7500cm}|p{4.350cm}}\hline\hline
\multicolumn{1}{c|}{神对约伯的心意}&\multicolumn{2}{c}{患難的态度}\\\cline{2-3}
&\multicolumn{1}{c|}{四个不可}&\multicolumn{1}{c}{不可忘記稱讚神為大}\\\hline
神藉著困苦救拔困苦人，趁他們受欺壓，開通他們的耳朵。神也必引你出離患難，進入寬闊不狹窄之地，擺在你席上的必滿有肥甘。
&但你滿口有惡人批評的言語，判斷和刑罰抓住你。\textcolor{red}{不可}容憤怒觸動你，使你不服責罰，也\textcolor{red}{不可}因贖價大就偏行。你的呼求(資財)，或是你一切的勢力，果有靈驗，叫你不受患難嗎？\textcolor{red}{不要}切慕黑夜，就是眾民在本處被除滅的時候。你要謹慎，\textcolor{red}{不可}重看罪孽，因你選擇罪孽過於選擇苦難。
&神行事有高大的能力，教訓人的有誰像他呢？誰派定他的道路？誰能說『你所行的不義』？你\textcolor{red}{不可}忘記稱讚他所行的為大，就是人所歌頌的。\\\hline
\end{tabular}
\end{table} 


\subsubsection{称颂神的作为 36:25-37:24}
\subsubsubsection{神在自然界的作为 36:25-30}
他所行的，萬人都看見，世人也從遠處觀看。 26 神為大，我們不能全知，他的年數不能測度。 27 他吸取水點，這水點從雲霧中就變成雨， 28 雲彩將雨落下，沛然降於世人。 29 誰能明白雲彩如何鋪張和神行宮的雷聲呢？ 30 他將亮光普照在自己的四圍，他又遮覆海底。
\subsubsubsection{神审判众民，甚至牲畜借着雷雨都知道神何时发怒 36:31-33}
他用這些審判眾民，且賜豐富的糧食。 32 他以電光遮手，命閃電擊中敵人（中了靶子）。 33 所發的雷聲顯明他的作為，又向牲畜指明要起暴風。
\subsubsubsection{动物都服从神，自然界随着神手而向人施行責罰，潤地，或施行慈愛 37:1-13}
因此我心戰兢，從原處移動。 2 聽啊，神轟轟的聲音，是他口中所發的響聲。 3 他發響聲震遍天下，發電光閃到地極。 4 隨後人聽見有雷聲轟轟，大發威嚴，雷電接連不斷。 5 神發出奇妙的雷聲，他行大事，我們不能測透。 6 他對雪說：『要降在地上』，對大雨和暴雨也是這樣說。 7 他封住各人的手，叫所造的萬人，都曉得他的作為。 8 百獸進入穴中，臥在洞內。 9 暴風出於南宮，寒冷出於北方。 10 神噓氣成冰，寬闊之水也都凝結。 11 他使密雲盛滿水氣，布散電光之雲。 12 這雲是藉他的指引遊行旋轉，得以在全地面上，行他一切所吩咐的， 13 或為責罰，或為潤地，或為施行慈愛。
\subsubsubsection{劝約伯留心聽，要站立思想神奇妙的作為 37:14-17}
約伯啊，你要留心聽，要站立思想神奇妙的作為。 15 神如何吩咐這些，如何使雲中的電光照耀，你知道嗎？ 16 雲彩如何浮於空中，那知識全備者奇妙的作為，你知道嗎？ 17 南風使地寂靜，你的衣服就如火熱，你知道嗎？
\subsubsubsection{人类是何等愚昧和渺小 37:18-24} 
\begin{table}[H]
\centering
\begin{tabular}{p{7.15cm}|p{9.7500cm}}\hline\hline
\multicolumn{1}{c|}{我們愚昧不知道該對他說什麼話}&\multicolumn{1}{c}{必不苦待人}\\\cline{1-2}
你豈能與神同鋪穹蒼嗎？這穹蒼堅硬，如同鑄成的鏡子。 19 我們愚昧不能陳說，請你指教我們該對他說什麼話。 20 人豈可說『我願與他說話』？豈有人自願滅亡嗎？
&現在有雲遮蔽，人不得見穹蒼的光亮；但風吹過，天又發晴。 22 金光出於北方，在神那裡有可怕的威嚴。 23 論到全能者，我們不能測度，他大有能力，有公平和大義，必不苦待人。 24 所以人敬畏他，凡自以為心中有智慧的人，他都不顧念。\\\hline
\end{tabular}
\end{table} 


\section{神的安慰 38-42}
\subsection{神三问約伯 38-41}
\subsubsection{以大自然的创造一问約伯 38:1-38}
\begin{itemize}
\itemsep-.45em
\item 神從旋風中回答約伯 38:1-3\\
那時，耶和華從旋風中回答約伯說： 2 「誰用無知的言語，使我的旨意暗昧不明？ 3 你要如勇士束腰，我問你，你可以指示我。
\item 神立大地根基 38:4-7\\
我立大地根基的時候，你在哪裡呢？你若有聰明，只管說吧。 5 你若曉得就說，是誰定地的尺度？是誰把準繩拉在其上？ 6 地的根基安置在何處？地的角石是誰安放的？ 7 那時晨星一同歌唱，神的眾子也都歡呼。
\item 為海水定界限 38:8-11\\
海水衝出，如出胎胞，那時誰將它關閉呢？ 9 是我用雲彩當海的衣服，用幽暗當包裹它的布， 10 為它定界限，又安門和閂， 11 說：『你只可到這裡，不可越過，你狂傲的浪要到此止住。』
\item 命定晨光 38:12-15\\
你自生以來，曾命定晨光，使清晨的日光知道本位， 13 叫這光普照地的四極，將惡人從其中驅逐出來嗎？ 14 因這光，地面改變如泥上印印，萬物出現如衣服一樣。 15 亮光不照惡人，強橫的膀臂也必折斷。
\item 造黑暗和光明 38:16-21 \\
你曾進到海源，或在深淵的隱密處行走嗎？ 17 死亡的門曾向你顯露嗎？死蔭的門你曾見過嗎？ 18 地的廣大你能明透嗎？你若全知道，只管說吧。光明的居所從何而至？黑暗的本位在於何處？ 20 你能帶到本境，能看明其室之路嗎？ 21 你總知道，因為你早已生在世上，你日子的數目也多。 
\item 雪庫,雹倉,雷,電,雨,雲,眾星受造各有其目的 38:22-37\\
你曾進入雪庫，或見過雹倉嗎？ 23 這雪雹乃是我為降災並打仗和爭戰的日子所預備的。 24 光亮從何路分開？東風從何路分散遍地？

25 「誰為雨水分道？誰為雷電開路？ 26 使雨降在無人之地，無人居住的曠野， 27 使荒廢淒涼之地得以豐足，青草得以發生。 28 雨有父嗎？露水珠是誰生的呢？ 29 冰出於誰的胎？天上的霜是誰生的呢？ 30 諸水堅硬(隱藏)如石頭，深淵之面凝結成冰。

31 「你能繫住昴星的結嗎？能解開參星的帶嗎？ 32 你能按時領出十二宮嗎？能引導北斗和隨它的眾星(子)嗎？ 33 你知道天的定例嗎？能使地歸在天的權下嗎？

34 「你能向雲彩揚起聲來，使傾盆的雨遮蓋你嗎？ 35 你能發出閃電，叫它行去，使它對你說『我們在這裡』？ 36 誰將智慧放在懷中？誰將聰明賜於心內？ 37 誰能用智慧數算雲彩呢？塵土聚集成團，土塊緊緊結連， 38 那時，誰能傾倒天上的瓶呢？
\end{itemize}

\subsubsection{以自然界的动物二问约伯 38:39-39:40}
\begin{itemize}
\itemsep-.45em
\item 獅子抓取食物,烏鴉哀告神 38:39-41\\
母獅子在洞中蹲伏，少壯獅子在隱密處埋伏， 40 你能為牠們抓取食物，使牠們飽足嗎？ 41 烏鴉之雛因無食物飛來飛去，哀告神，那時，誰為牠預備食物呢？

\item 野山羊,母鹿 39:1-4\\
山巖間的野山羊幾時生產，你知道嗎？母鹿下犢之期，你能察定嗎？ 2 牠們懷胎的月數，你能數算嗎？牠們幾時生產，你能曉得嗎？ 3 牠們屈身，將子生下，就除掉疼痛。 4 這子漸漸肥壯，在荒野長大，去而不回。

\item 野驢 39:5-8\\
誰放野驢出去自由？誰解開快驢的繩索？ 6 我使曠野做牠的住處，使鹹地當牠的居所。 7 牠嗤笑城內的喧嚷，不聽趕牲口的喝聲。 8 遍山是牠的草場，牠尋找各樣青綠之物。

\item 野牛 39:9-12\\
野牛豈肯服侍你？豈肯住在你的槽旁？ 10 你豈能用套繩將野牛籠在犁溝之間？牠豈肯隨你耙山谷之地？ 11 豈可因牠的力大就倚靠牠？豈可把你的工交給牠做嗎？ 12 豈可信靠牠把你的糧食運到家，又收聚你禾場上的穀嗎？

\item 鴕鳥 39:13-18\\
鴕鳥的翅膀歡然搧展，豈是顯慈愛的翎毛和羽毛嗎？ 14 因牠把蛋留在地上，在塵土中使得溫暖， 15 卻想不到被腳踹碎，或被野獸踐踏。 16 牠忍心待雛，似乎不是自己的；雖然徒受勞苦，也不為雛懼怕。 17 因為神使牠沒有智慧，也未將悟性賜給牠。 18 牠幾時挺身展開翅膀，就嗤笑馬和騎馬的人。

\item 馬 39:19-25 \\
馬的大力是你所賜的嗎？牠頸項上挓挲的鬃是你給牠披上的嗎？ 20 是你叫牠跳躍像蝗蟲嗎？牠噴氣之威使人驚惶。 21 牠在谷中刨地，自喜其力，牠出去迎接佩帶兵器的人。 22 牠嗤笑可怕的事並不驚惶，也不因刀劍退回。 23 箭袋和發亮的槍並短槍，在牠身上錚錚有聲。 24 牠發猛烈的怒氣將地吞下，一聽角聲就不耐站立。 25 角每發聲，牠說『呵哈』，牠從遠處聞著戰氣，又聽見軍長大發雷聲和兵丁呐喊。

\item 鷹雀 39:26-30\\
鷹雀飛翔，展開翅膀一直向南，豈是藉你的智慧嗎？ 27 大鷹上騰在高處搭窩，豈是聽你的吩咐嗎？ 28 牠住在山巖，以山峰和堅固之所為家， 29 從那裡窺看食物，眼睛遠遠觀望。 30 牠的雛也咂血；被殺的人在哪裡，牠也在哪裡。
\end{itemize}

\subsubsection{约伯的第一次回答-向神自卑 40:1-5}
耶和華又對約伯說： 2 「強辯的豈可與全能者爭論嗎？與神辯駁的可以回答這些吧！」3 於是約伯回答耶和華說： 4 「我是卑賤的，我用什麼回答你呢？只好用手摀口。 5 我說了一次，再不回答；說了兩次，就不再說。」


\subsubsection{神以能力的事物三问约伯 40:6-41:34}
\begin{itemize}
\itemsep-.45em
\item 制伏驕傲的人，惡人踐踏在本處 40:6-14\\
於是耶和華從旋風中回答約伯說： 7 「你要如勇士束腰，我問你，你可以指示我。
你豈可廢棄我所擬定的？豈可定我有罪，好顯自己為義嗎？ 9 你有神那樣的膀臂嗎？你能像他發雷聲嗎？
10 「你要以榮耀莊嚴為裝飾，以尊榮威嚴為衣服。 11 要發出你滿溢的怒氣，見一切驕傲的人，使他降卑， 12 見一切驕傲的人，將他制伏，把惡人踐踏在本處。 13 將他們一同隱藏在塵土中，把他們的臉蒙蔽在隱密處。 14 我就認你右手能以救自己。

\item 河馬的能力 40:15-24\\
你且觀看河馬，我造你也造牠，牠吃草與牛一樣。 16 牠的氣力在腰間，能力在肚腹的筋上。 17 牠搖動尾巴如香柏樹，牠大腿的筋互相聯絡。 18 牠的骨頭好像銅管，牠的肢體彷彿鐵棍。 19 牠在神所造的物中為首，創造牠的給牠刀劍。 20 諸山給牠出食物，也是百獸遊玩之處。 21 牠伏在蓮葉之下，臥在蘆葦隱密處和水窪子裡。 22 蓮葉的陰涼遮蔽牠，溪旁的柳樹環繞牠。 23 河水氾濫，牠不發戰，就是約旦河的水漲到牠口邊，也是安然。 24 在牠防備的時候，誰能捉拿牠？誰能牢籠牠，穿牠的鼻子呢？

\item 鱷魚的肢體和其大力 41:1-34\\
你能用魚鉤釣上鱷魚嗎？能用繩子壓下牠的舌頭嗎？ 2 你能用繩索穿牠的鼻子嗎？能用鉤穿牠的腮骨嗎？ 3 牠豈向你連連懇求，說柔和的話嗎？ 4 豈肯與你立約，使你拿牠永遠做奴僕嗎？ 5 你豈可拿牠當雀鳥玩耍嗎？豈可為你的幼女將牠拴住嗎？ 6 搭夥的漁夫豈可拿牠當貨物嗎？能把牠分給商人嗎？ 7 你能用倒鉤槍扎滿牠的皮，能用魚叉叉滿牠的頭嗎？ 8 你按手在牠身上，想與牠爭戰，就不再這樣行吧。 9 人指望捉拿牠是徒然的，一見牠，豈不喪膽嗎？ 10 沒有那麼凶猛的人敢惹牠。這樣，誰能在我面前站立得住呢？ 11 誰先給我什麼，使我償還呢？天下萬物都是我的。

12 「論到鱷魚的肢體和其大力，並美好的骨骼，我不能緘默不言。 13 誰能剝牠的外衣？誰能進牠上下牙骨之間呢？ 14 誰能開牠的腮頰？牠牙齒四圍是可畏的。 15 牠以堅固的鱗甲為可誇，緊緊合閉，封得嚴密。 16 這鱗甲一一相連，甚至氣不得透入其間， 17 都是互相聯絡，膠結不能分離。 18 牠打噴嚏就發出光來，牠眼睛好像早晨的光線 (眼皮)。 19 從牠口中發出燒著的火把，與飛迸的火星。 20 從牠鼻孔冒出煙來，如燒開的鍋和點著的蘆葦。 21 牠的氣點著煤炭，有火焰從牠口中發出。 22 牠頸項中存著勁力，在牠面前的都恐嚇蹦跳。 23 牠的肉塊互相聯絡，緊貼其身，不能搖動。 24 牠的心結實如石頭，如下磨石那樣結實。 25 牠一起來，勇士都驚恐，心裡慌亂，便都昏迷。 26 人若用刀，用槍，用標槍，用尖槍扎牠，都是無用。 27 牠以鐵為乾草，以銅為爛木。 28 箭不能恐嚇牠使牠逃避，彈石在牠看為碎秸。 29 棍棒算為禾秸，牠嗤笑短槍颼的響聲。 30 牠肚腹下如尖瓦片，牠如釘耙經過淤泥。 31 牠使深淵開滾如鍋，使洋海如鍋中的膏油。 32 牠行的路隨後發光，令人想深淵如同白髮。 33 在地上沒有像牠造的那樣，無所懼怕。 34 凡高大的，牠無不藐視，牠在驕傲的水族上做王。」
\end{itemize}

\subsubsection{約伯的第二次回答-認罪自責 42:1-6}
約伯回答耶和華說： 2 「我知道，你萬事都能做，你的旨意不能攔阻。 3 誰用無知的言語使你的旨意隱藏呢？我所說的，是我不明白的；這些事太奇妙，是我不知道的。 4 求你聽我，我要說話；我問你，求你指示我。 5 我從前風聞有你，現在親眼看見你。 6 因此我厭惡自己(我的言語)，在塵土和爐灰中懊悔。」

\subsection{耶和華責其三友 42:7-9}
耶和華對約伯說話以後，就對提幔人以利法說：「我的怒氣向你和你兩個朋友發作，因為你們議論我，不如我的僕人約伯說的是。 8 現在你們要取七隻公牛、七隻公羊，到我僕人約伯那裡去，為自己獻上燔祭，我的僕人約伯就為你們祈禱。我因悅納他，就不按你們的愚妄辦你們。你們議論我，不如我的僕人約伯說的是。」 9 於是提幔人以利法、書亞人比勒達、拿瑪人瑣法照著耶和華所吩咐的去行，耶和華就悅納約伯。


\subsection{神賜福約伯較昔加倍 42:10-17}
約伯為他的朋友祈禱，耶和華就使約伯從苦境(擄掠)轉回，並且耶和華賜給他的，比他從前所有的加倍。 11 約伯的弟兄姐妹和以先所認識的人都來見他，在他家裡一同吃飯，又論到耶和華所降於他的一切災禍，都為他悲傷安慰他，每人也送他一塊銀子和一個金環。 12 這樣，耶和華後來賜福給約伯比先前更多。他有一萬四千羊，六千駱駝，一千對牛，一千母驢。 13 他也有七個兒子，三個女兒。 14 他給長女起名叫耶米瑪，次女叫基洗亞，三女叫基連哈樸。 15 在那全地的婦女中，找不著像約伯的女兒那樣美貌。她們的父親使她們在弟兄中得產業。此後，約伯又活了一百四十年，得見他的兒孫，直到四代。 17 這樣，約伯年紀老邁，日子滿足而死。







